% =============================================================================
% GRADUATION PROJECT THESIS — MASTER FORMATTING FILE v2
% =============================================================================
% Based on cross-research of: IEEE, MIT, Stanford, Oxford, Cambridge,
% APA 7th, Chicago/Turabian, ACM, Springer svmono, Cairo University FCI,
% Ain Shams Engineering, and Egyptian Supreme Council of Universities.
%
% PRIMARY STANDARD: FCI Cairo University (CS graduation project)
% SECONDARY: MIT/Oxford/Cambridge best practices
% TABLE STYLE:   APA 7th (3 horizontal rules only, no vertical rules)
% REFERENCES:    IEEE style (standard for CS departments)
%
% Compile with: XeLaTeX or LuaLaTeX
% Usage: % =============================================================================
% GRADUATION PROJECT THESIS — MASTER FORMATTING FILE v2
% =============================================================================
% Based on cross-research of: IEEE, MIT, Stanford, Oxford, Cambridge,
% APA 7th, Chicago/Turabian, ACM, Springer svmono, Cairo University FCI,
% Ain Shams Engineering, and Egyptian Supreme Council of Universities.
%
% PRIMARY STANDARD: FCI Cairo University (CS graduation project)
% SECONDARY: MIT/Oxford/Cambridge best practices
% TABLE STYLE:   APA 7th (3 horizontal rules only, no vertical rules)
% REFERENCES:    IEEE style (standard for CS departments)
%
% Compile with: XeLaTeX or LuaLaTeX
% Usage: % =============================================================================
% GRADUATION PROJECT THESIS — MASTER FORMATTING FILE v2
% =============================================================================
% Based on cross-research of: IEEE, MIT, Stanford, Oxford, Cambridge,
% APA 7th, Chicago/Turabian, ACM, Springer svmono, Cairo University FCI,
% Ain Shams Engineering, and Egyptian Supreme Council of Universities.
%
% PRIMARY STANDARD: FCI Cairo University (CS graduation project)
% SECONDARY: MIT/Oxford/Cambridge best practices
% TABLE STYLE:   APA 7th (3 horizontal rules only, no vertical rules)
% REFERENCES:    IEEE style (standard for CS departments)
%
% Compile with: XeLaTeX or LuaLaTeX
% Usage: % =============================================================================
% GRADUATION PROJECT THESIS — MASTER FORMATTING FILE v2
% =============================================================================
% Based on cross-research of: IEEE, MIT, Stanford, Oxford, Cambridge,
% APA 7th, Chicago/Turabian, ACM, Springer svmono, Cairo University FCI,
% Ain Shams Engineering, and Egyptian Supreme Council of Universities.
%
% PRIMARY STANDARD: FCI Cairo University (CS graduation project)
% SECONDARY: MIT/Oxford/Cambridge best practices
% TABLE STYLE:   APA 7th (3 horizontal rules only, no vertical rules)
% REFERENCES:    IEEE style (standard for CS departments)
%
% Compile with: XeLaTeX or LuaLaTeX
% Usage: \input{thesis_formatting} in main.tex preamble
% =============================================================================

% ─── DOCUMENT CLASS ──────────────────────────────────────────────────────────
% book class: best for multi-chapter theses (front/main/back matter support)
% 12pt: universal standard (Oxford, Cambridge, APA, Chicago, Cairo Univ)
% a4paper: international standard (Oxford, Cambridge, Egyptian universities)
% oneside: standard for single-sided printing
% openright: chapters start on right-hand (odd) pages — MIT/Oxford standard
\documentclass[12pt, a4paper, oneside, openright]{report}
% For two-sided printing (required by some universities), change to:
% \documentclass[12pt, a4paper, twoside, openright]{report}

% ─── USER CONFIGURATION ─────────────────────────────────────────────────────
\newcommand{\UniversityNameEN}{Faculty of Computers and Information}
\newcommand{\UniversityNameAR}{كلية الحاسبات والمعلومات}
\newcommand{\DepartmentNameEN}{Computer Science Department}
\newcommand{\DepartmentNameAR}{قسم علوم الحاسب}
\newcommand{\UniversityFullEN}{Cairo University}
\newcommand{\ThesisTitle}{Quiz Generation via Knowledge Distillation and Fine-Tuning}
\newcommand{\AuthorName}{Rehab Hamdy}
\newcommand{\SupervisorName}{Prof.\ XXX}
\newcommand{\AcademicYear}{2025--2026}

% ─── PAGE GEOMETRY ───────────────────────────────────────────────────────────
% FCI Cairo: 3cm left (binding), 2.5cm others
% Egyptian Supreme Council: 4cm binding side, 2.5cm others
% MIT/Stanford: 1.5" (38mm) inner for binding
% Compromise: 3.5cm inner (binding), 2.5cm others — satisfies all standards
\usepackage[
  a4paper,
  top=2.5cm,
  bottom=2.5cm,
  inner=3.5cm,                          % binding side (left for English)
  outer=2.5cm,
  headheight=14.5pt,
  headsep=15mm,
  footskip=20mm,
  marginparwidth=0pt,
  marginparsep=0pt
]{geometry}

% ─── LANGUAGE & ENCODING ─────────────────────────────────────────────────────
\usepackage[utf8]{inputenc}              % UTF-8 input
\usepackage[T1]{fontenc}                 % T1 font encoding
\usepackage[english]{babel}              % English hyphenation
% Add arabic option if bilingual content: \usepackage[english,arabic]{babel}

% ─── FONTS ───────────────────────────────────────────────────────────────────
% Times New Roman: universal thesis standard (Oxford, Cambridge, APA, Chicago,
% MIT, Stanford, FCI Cairo, Ain Shams)
\usepackage{fontspec}                    % XeLaTeX / LuaLaTeX font selection
\setmainfont{Times New Roman}            % Body serif font — 12pt standard
\setsansfont{Arial}                      % Sans-serif for captions, headers
\setmonofont{Courier New}[Scale=0.90]    % Monospace for code — 10~11pt equiv.

% Fallback if TNR not available:
% \setmainfont{TeX Gyre Termes}         % Free Times clone
% \setsansfont{TeX Gyre Heros}          % Free Helvetica/Arial clone
% \setmonofont{TeX Gyre Cursor}         % Free Courier clone

% ─── LINE SPACING ────────────────────────────────────────────────────────────
% FCI Cairo: 1.5 line spacing
% Oxford/Cambridge: 1.5 or double
% APA: double
% MIT/Stanford: 1.5 or double
% → 1.5 is the most common for FCI; use \doublespacing if your dept. requires
\usepackage{setspace}
\onehalfspacing

% ─── PARAGRAPH FORMATTING ────────────────────────────────────────────────────
% APA / Chicago / FCI: 0.5" (12.7mm) first-line indent
% APA: NO extra space between paragraphs
\usepackage{indentfirst}                 % indent first paragraph after heading
\setlength{\parindent}{12.7mm}           % 0.5 inch first-line indent
\setlength{\parskip}{0pt}                % no extra vertical space between paras

% ─── SECTIONAL HEADINGS ─────────────────────────────────────────────────────
% Heading hierarchy based on APA 7th + FCI Cairo synthesis:
%   Chapter  → centered, bold, 16pt (FCI standard)
%   Section  → left-aligned, bold, 14pt (FCI standard)
%   Subsection → left-aligned, bold, 12pt (FCI standard)
%   Subsubsection → left-aligned, bold italic, 12pt
%   Paragraph → indented, bold, run-in, period, 12pt
\usepackage{titlesec}
\usepackage{titletoc}

% ---- Chapter (APA L1 equivalent) ----
% FCI: bold, centered, 16pt
% Oxford/Cambridge: centered, bold, 14-16pt
\titleformat{\chapter}[display]
  {\normalfont\centering\bfseries}
  {\MakeUppercase{\chaptertitlename}\ \thechapter}
  {20pt}
  {\fontsize{16pt}{19.2pt}\selectfont\bfseries}
\titlespacing*{\chapter}{0pt}{-10pt}{30pt}

% ---- Section (APA L2 equivalent) ----
% FCI: bold, left-aligned, 14pt
% APA: Flush Left, Bold, Title Case
\titleformat{\section}
  {\normalfont\bfseries\fontsize{14pt}{16.8pt}\selectfont}
  {\thesection}
  {1em}
  {}
\titlespacing*{\section}{0pt}{18pt plus 4pt minus 2pt}{10pt plus 2pt minus 2pt}

% ---- Subsection (APA L3 equivalent) ----
% FCI: bold, left-aligned, 12pt
% APA: Flush Left, Bold Italic, Title Case
\titleformat{\subsection}
  {\normalfont\bfseries\fontsize{12pt}{14.4pt}\selectfont}
  {\thesubsection}
  {1em}
  {}
\titlespacing*{\subsection}{0pt}{14pt plus 4pt minus 2pt}{8pt plus 2pt minus 2pt}

% ---- Subsubsection (APA L4 equivalent) ----
% APA: Indented, Bold, Title Case, Period., Text continues
% Chicago: Flush Left, Italic
\titleformat{\subsubsection}
  {\normalfont\bfseries\itshape}
  {\thesubsubsection}
  {1em}
  {}
\titlespacing*{\subsubsection}{0pt}{12pt plus 4pt minus 2pt}{6pt plus 2pt minus 2pt}

% ---- Paragraph (APA L5 equivalent) ----
% APA: Indented, Bold Italic, Title Case, Period., Text continues
\titleformat{\paragraph}[runin]
  {\normalfont\bfseries\itshape}
  {\theparagraph}
  {1em}
  {}
\titlespacing*{\paragraph}{\parindent}{10pt plus 4pt minus 2pt}{1em plus 2pt minus 2pt}

% ─── TABLE OF CONTENTS FORMATTING ───────────────────────────────────────────
% Oxford-style: clean, professional dot leaders
\titlecontents{chapter}
  [0em]
  {\bfseries\vspace{6pt}}
  {\contentslabel{2em}}
  {\hspace*{2em}}
  {\titlerule*[8pt]{$\cdot$}\contentspage}

\titlecontents{section}
  [2em]
  {}
  {\contentslabel{2.5em}}
  {\hspace*{2.5em}}
  {\titlerule*[8pt]{$\cdot$}\contentspage}

\titlecontents{subsection}
  [4.5em]
  {}
  {\contentslabel{3em}}
  {\hspace*{3em}}
  {\titlerule*[8pt]{$\cdot$}\contentspage}

\titlecontents{subsubsection}
  [7.5em]
  {\small}
  {\contentslabel{3.5em}}
  {\hspace*{3.5em}}
  {\titlerule*[8pt]{$\cdot$}\contentspage}

% ─── NUMBERING DEPTH ────────────────────────────────────────────────────────
\setcounter{secnumdepth}{4}              % number up to subsubsection
\setcounter{tocdepth}{3}                 % show up to subsection in TOC
                                          % (change to 2 for sections only)

% ─── FIGURE & TABLE NUMBERING ───────────────────────────────────────────────
% Per-chapter numbering: Fig 3.1, Table 3.2 (MIT, Oxford, FCI standard)
\usepackage{chngcntr}
\counterwithin{figure}{chapter}
\counterwithin{table}{chapter}
% For continuous numbering (Fig 1, Fig 2, ...), use instead:
% \counterwithout{figure}{chapter}
% \counterwithout{table}{chapter}

% ─── PACKAGES — TABLES ──────────────────────────────────────────────────────
% APA 7th: ONLY 3 horizontal rules — top, header-body separator, bottom
% NO vertical rules. This is the universal academic standard.
\usepackage{booktabs}                    % \toprule, \midrule, \bottomrule
\usepackage{tabularx}                    % flexible-width tables
\usepackage{longtable}                   % tables spanning multiple pages
\usepackage{multirow}                    % merged cells vertically
\usepackage{makecell}                    % line breaks inside cells
\usepackage{array}                       % extended column types
\usepackage{colortbl}                    % colored cells (use sparingly)

% ─── PACKAGES — FIGURES ─────────────────────────────────────────────────────
\usepackage{graphicx}                    % \includegraphics
\usepackage{float}                       % [H] placement specifier
\usepackage{subcaption}                  % sub-figures (a), (b), (c)
\usepackage{wrapfig}                     % text wrapping around figures
\usepackage{tikz}
\usetikzlibrary{
  shapes, arrows, positioning, calc,
  fit, backgrounds, decorations.pathreplacing,
  matrix, shadows, patterns
}

% ─── PACKAGES — CAPTIONS ────────────────────────────────────────────────────
% Universal standard: Table caption ABOVE, Figure caption BELOW
% APA: Table title italic, "Note." below table
\usepackage[
  font={small},
  labelfont={bf},
  labelsep=quad,
  justification=centerlast,
  skip=8pt
]{caption}

% Figure caption: below figure (ALL standards agree)
\captionsetup[figure]{
  font={small},
  labelfont={bf},
  labelsep=quad,
  justification=centerlast,
  skip=6pt,
  position=bottom
}

% Table caption: above table (ALL standards agree)
\captionsetup[table]{
  font={small},
  labelfont={bf},
  labelsep=quad,
  justification=centerlast,
  skip=6pt,
  position=top
}

% Sub-figure caption
\captionsetup[subfigure]{
  font={footnote},
  labelfont={bf},
  labelsep=space,
  justification=centerlast,
  skip=4pt
}

% ─── PACKAGES — HEADERS & FOOTERS ───────────────────────────────────────────
\usepackage{fancyhdr}

% ─── PACKAGES — COLORS ──────────────────────────────────────────────────────
\usepackage{xcolor}
\usepackage[table]{xcolor}

% ─── PACKAGES — MATHEMATICS ─────────────────────────────────────────────────
\usepackage{amsmath}
\usepackage{amssymb}
\usepackage{amsfonts}
\usepackage{mathtools}
\usepackage{bm}

% ─── PACKAGES — ALGORITHMS ──────────────────────────────────────────────────
\usepackage{algorithm}
\usepackage{algpseudocode}
\usepackage{algorithmic}

% ─── PACKAGES — CODE LISTINGS ───────────────────────────────────────────────
\usepackage{listings}

% ─── PACKAGES — HYPERLINKS ──────────────────────────────────────────────────
\usepackage[
  colorlinks=true,
  linkcolor=ThesisBlue,
  citecolor=ThesisGreen,
  urlcolor=ThesisNavy,
  bookmarks=true,
  bookmarksopen=true,
  bookmarksnumbered=true,
  pdfauthor={\AuthorName},
  pdftitle={\ThesisTitle}
]{hyperref}

% ─── PACKAGES — BIBLIOGRAPHY ────────────────────────────────────────────────
% IEEE style: standard for CS departments (FCI, Ain Shams, MIT EECS)
\usepackage[backend=biber, style=ieee]{biblatex}
% \addbibresource{references.bib}        % uncomment and set your .bib file

% ─── PACKAGES — GLOSSARY / ACRONYMS ─────────────────────────────────────────
\usepackage[acronym, toc]{glossaries}
% \makeglossaries

% ─── PACKAGES — MISCELLANEOUS ───────────────────────────────────────────────
\usepackage{enumitem}                    % customized lists
\usepackage{pifont}                      % ding symbols
\usepackage{calc}
\usepackage{ifthen}
\usepackage{afterpage}
\usepackage{emptypage}                   % blank pages have no headers/footers
\usepackage{appendix}
\usepackage{pdflscape}                   % landscape pages
\usepackage{placeins}                    % \FloatBarrier
\usepackage{siunitx}                     % SI units
\usepackage{tcolorbox}                   % colored boxes for callouts
\usepackage{csquotes}                    % proper quoting (required by biblatex)

% ─── COLOR PALETTE ──────────────────────────────────────────────────────────
% Muted academic colors — professional, not distracting
% Inspired by Oxford/Cambridge thesis templates

% ---- Primary ----
\definecolor{ThesisBlue}{RGB}{0, 70, 140}         % headings, links
\definecolor{ThesisNavy}{RGB}{0, 40, 85}           % deep emphasis
\definecolor{ThesisGreen}{RGB}{0, 120, 60}         % citations, positive
\definecolor{ThesisRed}{RGB}{180, 30, 30}           % warnings, negative
\definecolor{ThesisOrange}{RGB}{210, 105, 0}        % highlights

% ---- Table ----
\definecolor{TableHeader}{RGB}{0, 70, 140}          % header row background
\definecolor{TableHeaderText}{RGB}{255, 255, 255}   % header row text
\definecolor{TableRowA}{RGB}{240, 245, 250}         % alternating row A
\definecolor{TableRowB}{RGB}{255, 255, 255}         % alternating row B

% ---- Code ----
\definecolor{CodeBG}{RGB}{248, 248, 248}
\definecolor{CodeFrame}{RGB}{200, 200, 200}
\definecolor{CodeKeyword}{RGB}{0, 0, 180}
\definecolor{CodeString}{RGB}{0, 128, 0}
\definecolor{CodeComment}{RGB}{128, 128, 128}
\definecolor{CodeNumber}{RGB}{160, 100, 0}

% ---- Callout Boxes ----
\definecolor{NoteBoxBG}{RGB}{232, 244, 255}
\definecolor{NoteBoxFrame}{RGB}{0, 70, 140}
\definecolor{WarningBoxBG}{RGB}{255, 243, 224}
\definecolor{WarningBoxFrame}{RGB}{210, 105, 0}
\definecolor{ImportantBoxBG}{RGB}{255, 232, 232}
\definecolor{ImportantBoxFrame}{RGB}{180, 30, 30}

% ─── PAGE STYLE — HEADERS & FOOTERS ────────────────────────────────────────
% Oxford/Cambridge: chapter title left header, page number right
% APA: page number top right (student paper) or running head + page# (pro)
% FCI: page number bottom center is common
% → Using Oxford-style with page number in header

\pagestyle{fancy}
\fancyhf{}

% ---- Header ----
\fancyhead[LE,RO]{\small\thepage}                     % page number
\fancyhead[RE]{\small\itshape\nouppercase{\leftmark}}  % chapter title (even)
\fancyhead[LO]{}                                       % leave blank (odd)
\renewcommand{\headrulewidth}{0.4pt}

% ---- Footer ----
\fancyfoot[C]{}
\fancyfoot[L]{\tiny\textcolor{gray}{\UniversityNameEN}}

% ---- Chapter-opening pages: no header, page number bottom center ----
\fancypagestyle{plain}{%
  \fancyhf{}%
  \fancyfoot[C]{\small\thepage}%
  \renewcommand{\headrulewidth}{0pt}%
}

% ─── TABLE FORMATTING ───────────────────────────────────────────────────────
% APA 7th (universally accepted):
%   - ONLY 3 horizontal rules: \toprule, \midrule, \bottomrule
%   - NO vertical rules
%   - Caption ABOVE table
%   - "Note." line below table (optional)
% Table text may be smaller (10pt) than body (12pt) — APA, Chicago

\renewcommand{\arraystretch}{1.30}       % row height (APA recommends 1.25-1.35)
\tabcolsep=8pt                           % column separation

% ---- Custom column types ----
\newcolumntype{L}[1]{>{\raggedright\arraybackslash}p{#1}}
\newcolumntype{C}[1]{>{\centering\arraybackslash}p{#1}}
\newcolumntype{R}[1]{>{\raggedleft\arraybackslash}p{#1}}
\newcolumntype{M}[1]{>{\centering\arraybackslash}m{#1}}

% ---- Styled table header command ----
% Use within booktabs tables for professional APA-style headers
\newcommand{\thead}[1]{\textbf{#1}}

% ---- Table Note command (APA style) ----
% Usage: \tablenote{Source: Author's calculations. n = 200.}
\newcommand{\tablenote}[1]{%
  \vspace{2pt}%
  {\small\textit{Note.}\quad #1}%
}

% ---- Professional table environment ----
% Usage:
%   \begin{proftable}{caption}{label}{col spec}
%     \toprule
%     \thead{Col 1} & \thead{Col 2} \\
%     \midrule
%     data & data \\
%     \bottomrule
%   \end{proftable}
\newenvironment{proftable}[3]{%
  \begin{table}[H]%
  \centering%
  \caption{#1}%
  \label{#2}%
  \small%
  \begin{tabular}{#3}%
}{%
  \end{tabular}%
  \end{table}%
}

% ---- Long table environment (multi-page) ----
% Same rules apply: booktabs only, no vertical rules
\newenvironment{proflongtable}[3]{%
  \begin{longtable}{#3}%
  \caption{#1}\label{#2}\\%
}{%
  \end{longtable}%
}

% ─── FIGURE PLACEMENT ───────────────────────────────────────────────────────
% APA: figures embedded near first mention, not floating far away
\renewcommand{\topfraction}{0.92}
\renewcommand{\bottomfraction}{0.85}
\renewcommand{\textfraction}{0.08}
\renewcommand{\floatpagefraction}{0.90}

% ─── LIST FORMATTING ────────────────────────────────────────────────────────
% APA: bullet lists allowed, no extra spacing between items
% Chicago: consistent indentation

\setitemize{
  leftmargin=2em,
  itemsep=4pt,
  parsep=2pt,
  topsep=6pt
}

\setenumerate{
  leftmargin=2em,
  itemsep=4pt,
  parsep=2pt,
  topsep=6pt,
  label=\arabic*.
}

\setdescription{
  leftmargin=2em,
  itemsep=4pt,
  parsep=2pt,
  topsep=6pt,
  font=\normalfont\bfseries
}

% ─── CODE LISTING FORMATTING ────────────────────────────────────────────────
% FCI Cairo: Courier New 10pt for code
\lstdefinestyle{ThesisCode}{
  backgroundcolor=\color{CodeBG},
  basicstyle=\ttfamily\footnotesize,
  keywordstyle=\color{CodeKeyword}\bfseries,
  stringstyle=\color{CodeString},
  commentstyle=\color{CodeComment}\itshape,
  numberstyle=\tiny\color{CodeNumber},
  numbers=left,
  numbersep=8pt,
  frame=single,
  rulecolor=\color{CodeFrame},
  framesep=5pt,
  xleftmargin=15pt,
  xrightmargin=5pt,
  breaklines=true,
  showstringspaces=false,
  tabsize=4,
  captionpos=b,
  columns=flexible,
  keepspaces=true,
  escapeinside={(*@}{@*)}
}
\lstset{style=ThesisCode}

\lstdefinestyle{PythonCode}{style=ThesisCode, language=Python}
\lstdefinestyle{BashCode}{style=ThesisCode, language=bash}

\newcommand{\code}[1]{\texttt{\small #1}}

% ─── ALGORITHM / PSEUDOCODE FORMATTING ──────────────────────────────────────
\algrenewcommand\algorithmicrequire{\textbf{Input:}}
\algrenewcommand\algorithmicensure{\textbf{Output:}}
\floatname{algorithm}{Algorithm}
\renewcommand{\algorithmiccomment}[1]{\hfill$\triangleright$ \textit{#1}}

% ─── MATHEMATICS FORMATTING ─────────────────────────────────────────────────
% Per-chapter equation numbering: (3.1), (3.2) — MIT, Oxford, Springer standard
\DeclareMathOperator*{\argmin}{arg\,min}
\DeclareMathOperator*{\argmax}{arg\,max}
\DeclareMathOperator{\KL}{KL}
\DeclareMathOperator{\softmax}{softmax}
\DeclareMathOperator{\sigmoid}{sigmoid}
\DeclareMathOperator{\ReLU}{ReLU}

% Math shortcuts
\newcommand{\R}{\mathbb{R}}
\newcommand{\E}{\mathbb{E}}
\newcommand{\norm}[1]{\left\lVert #1 \right\rVert}
\newcommand{\abs}[1]{\left\lvert #1 \right\rvert}
\newcommand{\inner}[2]{\left\langle #1,\, #2 \right\rangle}
\newcommand{\condbar}{\,\vert\,}
\newcommand{\bigO}{\mathcal{O}}

\numberwithin{equation}{chapter}

% ─── CALLOUT BOXES ──────────────────────────────────────────────────────────
% Use sparingly — academic theses prefer inline text over boxes
% But useful for: definitions, key findings, warnings

\newtcolorbox{notebox}[1][Note]{%
  colback=NoteBoxBG,
  colframe=NoteBoxFrame,
  fonttitle=\bfseries\sffamily,
  title=#1,
  boxrule=0.8pt,
  arc=2pt,
  left=8pt, right=8pt, top=4pt, bottom=4pt,
  before skip=10pt, after skip=10pt
}

\newtcolorbox{warningbox}[1][Warning]{%
  colback=WarningBoxBG,
  colframe=WarningBoxFrame,
  fonttitle=\bfseries\sffamily,
  title=#1,
  boxrule=0.8pt,
  arc=2pt,
  left=8pt, right=8pt, top=4pt, bottom=4pt,
  before skip=10pt, after skip=10pt
}

\newtcolorbox{importantbox}[1][Important]{%
  colback=ImportantBoxBG,
  colframe=ImportantBoxFrame,
  fonttitle=\bfseries\sffamily,
  title=#1,
  boxrule=0.8pt,
  arc=2pt,
  left=8pt, right=8pt, top=4pt, bottom=4pt,
  before skip=10pt, after skip=10pt
}

\newtcolorbox{definitionbox}[1][Definition]{%
  colback=NoteBoxBG,
  colframe=ThesisBlue,
  fonttitle=\bfseries\sffamily\color{white},
  coltitle=white,
  colbacktitle=ThesisBlue,
  title=#1,
  boxrule=0.8pt,
  arc=2pt,
  left=8pt, right=8pt, top=4pt, bottom=4pt,
  before skip=10pt, after skip=10pt
}

% ─── CROSS-REFERENCE COMMANDS ───────────────────────────────────────────────
% Friendly references — universal academic style
\newcommand{\reffig}[1]{Figure~\ref{#1}}
\newcommand{\reftab}[1]{Table~\ref{#1}}
\newcommand{\refeq}[1]{Equation~\ref{#1}}
\newcommand{\refch}[1]{Chapter~\ref{#1}}
\newcommand{\refsec}[1]{Section~\ref{#1}}
\newcommand{\refalg}[1]{Algorithm~\ref{#1}}
\newcommand{\refapp}[1]{Appendix~\ref{#1}}
\newcommand{\reflst}[1]{Listing~\ref{#1}}
\newcommand{\refpage}[1]{page~\pageref{#1}}

% ─── UTILITY COMMANDS ───────────────────────────────────────────────────────
\newcommand{\pct}[1]{#1\%}
\newcommand{\metric}[2]{\textbf{#1}\,=\,#2}

% ─── COVER PAGE ─────────────────────────────────────────────────────────────
% FCI Cairo style: University name, Faculty, Department, Title, Author, Supervisor
\newcommand{\MakeCoverPage}{%
  \begin{titlepage}
    \centering
    \vspace*{0.5cm}

    {\scshape\LARGE \UniversityFullEN \par}
    \vspace{0.3cm}
    {\scshape\Large \UniversityNameEN \par}
    \vspace{0.2cm}
    {\scshape\large \DepartmentNameEN \par}

    \vspace{2cm}

    {\fontsize{18pt}{21.6pt}\selectfont\bfseries \ThesisTitle \par}

    \vspace{1.5cm}

    {\Large Graduation Project Thesis\par}

    \vspace{2cm}

    \begin{minipage}{0.45\textwidth}
      \begin{flushleft}\large
        \textbf{Submitted by:}\\
        \AuthorName
      \end{flushleft}
    \end{minipage}
    \hfill
    \begin{minipage}{0.45\textwidth}
      \begin{flushright}\large
        \textbf{Supervised by:}\\
        \SupervisorName
      \end{flushright}
    \end{minipage}

    \vfill

    {\large \AcademicYear\par}

  \end{titlepage}
}

% ─── ABSTRACT PAGE ──────────────────────────────────────────────────────────
% FCI: Arabic abstract first, then English (or vice versa per department)
\newcommand{\MakeAbstractPage}[2]{%
  \chapter*{Abstract}
  \addcontentsline{toc}{chapter}{Abstract}
  #1
  \vfill
  \ifthenelse{\equal{#2}{}}{}{%
    \newpage
    \chapter*{الملخص}
    \addcontentsline{toc}{chapter}{الملخص}
    #2
  }
}

% ─── ACKNOWLEDGMENTS ────────────────────────────────────────────────────────
\newcommand{\MakeAcknowledgmentsPage}[1]{%
  \chapter*{Acknowledgments}
  \addcontentsline{toc}{chapter}{Acknowledgments}
  #1
}

% ─── DEDICATION ─────────────────────────────────────────────────────────────
\newcommand{\MakeDedicationPage}[1]{%
  \vspace*{\fill}
  \begin{center}
    \textit{\Large #1}
  \end{center}
  \vspace*{\fill}
}

% ─── APPENDIX FORMATTING ────────────────────────────────────────────────────
% APA: labeled Appendix A, Appendix B; each on new page
% FCI: similar structure
\renewcommand{\appendixname}{Appendix}

% ─── BIBLIOGRAPHY FORMATTING ────────────────────────────────────────────────
% IEEE style for CS departments (FCI, Ain Shams, MIT EECS)
\defbibheading{bibliography}{%
  \chapter*{Bibliography}%
  \addcontentsline{toc}{chapter}{Bibliography}%
}

% ─── FRONT MATTER / MAIN MATTER ─────────────────────────────────────────────
% MIT/Oxford/Cambridge standard:
%   Front matter: Roman numerals (i, ii, iii, ...)
%   Main matter:  Arabic numerals (1, 2, 3, ...)
%   Chapter 1 starts at page 1
\makeatletter
\renewcommand\frontmatter{%
  \cleardoublepage
  \pagenumbering{roman}
}

\renewcommand\mainmatter{%
  \cleardoublepage
  \pagenumbering{arabic}
}

\renewcommand\backmatter{%
  \cleardoublepage
  \setcounter{chapter}{0}%
}
\makeatother

% ─── DRAFT MARKERS ──────────────────────────────────────────────────────────
\newif\ifdraft
\drafttrue                              % set to \draftfalse for final version

\ifdraft
  \newcommand{\todo}[1]{\textcolor{ThesisRed}{\bfseries[TODO: #1]}}
  \newcommand{\FIXME}[1]{\textcolor{ThesisRed}{\bfseries[FIXME: #1]}}
  \newcommand{\CITE}[1]{\textcolor{ThesisOrange}{\bfseries[CITE: #1]}}
\else
  \newcommand{\todo}[1]{}
  \newcommand{\FIXME}[1]{}
  \newcommand{\CITE}[1]{\cite{#1}}
\fi

% ─── ORPHAN / WIDOW PREVENTION ──────────────────────────────────────────────
% Chicago/Oxford: avoid orphan/widow lines
\clubpenalty=10000
\widowpenalty=10000
\emergencystretch=1em
\tolerance=200
\setlength{\emergencystretch}{3em}

% ─── HYPHENATION EXCEPTIONS ─────────────────────────────────────────────────
\hyphenation{
  fine-tun-ing
  know-ledge
  distil-la-tion
  eval-ua-tion
  pre-train-ed
  multi-ple
  parameter
  hyper-parameter
  archi-tec-ture
}

% ─── SI UNIT FORMATTING ─────────────────────────────────────────────────────
\sisetup{
  detect-all,
  group-separator={,},
  group-minimum-digits=4,
  per-mode=symbol
}

% =============================================================================
% FORMATTING DECISIONS LOG — Based on Cross-Source Research
% =============================================================================
%
% | Decision              | Chosen Value       | Sources Supporting It          |
% |-----------------------|--------------------|--------------------------------|
% | Page size             | A4                 | Oxford, Cambridge, Cairo, FCI  |
% | Inner margin          | 3.5cm              | FCI (3cm), MIT (1.5"), avg     |
% | Outer margin          | 2.5cm              | FCI, Egyptian Council          |
% | Body font             | TNR 12pt           | ALL sources agree              |
% | Line spacing          | 1.5                | FCI, Cambridge, Oxford         |
% | Paragraph indent      | 12.7mm (0.5")      | APA, Chicago, FCI              |
% | Chapter heading       | 16pt bold centered | FCI, Oxford, Cambridge         |
% | Section heading       | 14pt bold left     | FCI standard                   |
% | Subsection heading    | 12pt bold left     | FCI, APA L3                    |
% | Table rules           | 3 only (booktabs)  | APA 7th, ALL academic stds     |
% | Table caption         | Above table        | APA, IEEE, ACM, ALL sources    |
% | Figure caption        | Below figure       | APA, IEEE, ACM, ALL sources    |
% | Equation numbering    | Per chapter        | MIT, Oxford, Springer          |
% | Figure numbering      | Per chapter        | MIT, Oxford, FCI               |
% | Bibliography style    | IEEE               | FCI CS dept, ACM, MIT EECS     |
% | Front matter pages    | Roman numerals     | MIT, Oxford, Cambridge, FCI    |
% | Main matter pages     | Arabic numerals    | ALL sources agree              |
% | TOC depth             | 3 levels           | Oxford, Cambridge              |
% | Secnum depth          | 4 levels           | Standard for theses            |
%
% =============================================================================

% =============================================================================
% END OF FORMATTING FILE v2
% =============================================================================
 in main.tex preamble
% =============================================================================

% ─── DOCUMENT CLASS ──────────────────────────────────────────────────────────
% book class: best for multi-chapter theses (front/main/back matter support)
% 12pt: universal standard (Oxford, Cambridge, APA, Chicago, Cairo Univ)
% a4paper: international standard (Oxford, Cambridge, Egyptian universities)
% oneside: standard for single-sided printing
% openright: chapters start on right-hand (odd) pages — MIT/Oxford standard
\documentclass[12pt, a4paper, oneside, openright]{report}
% For two-sided printing (required by some universities), change to:
% \documentclass[12pt, a4paper, twoside, openright]{report}

% ─── USER CONFIGURATION ─────────────────────────────────────────────────────
\newcommand{\UniversityNameEN}{Faculty of Computers and Information}
\newcommand{\UniversityNameAR}{كلية الحاسبات والمعلومات}
\newcommand{\DepartmentNameEN}{Computer Science Department}
\newcommand{\DepartmentNameAR}{قسم علوم الحاسب}
\newcommand{\UniversityFullEN}{Cairo University}
\newcommand{\ThesisTitle}{Quiz Generation via Knowledge Distillation and Fine-Tuning}
\newcommand{\AuthorName}{Rehab Hamdy}
\newcommand{\SupervisorName}{Prof.\ XXX}
\newcommand{\AcademicYear}{2025--2026}

% ─── PAGE GEOMETRY ───────────────────────────────────────────────────────────
% FCI Cairo: 3cm left (binding), 2.5cm others
% Egyptian Supreme Council: 4cm binding side, 2.5cm others
% MIT/Stanford: 1.5" (38mm) inner for binding
% Compromise: 3.5cm inner (binding), 2.5cm others — satisfies all standards
\usepackage[
  a4paper,
  top=2.5cm,
  bottom=2.5cm,
  inner=3.5cm,                          % binding side (left for English)
  outer=2.5cm,
  headheight=14.5pt,
  headsep=15mm,
  footskip=20mm,
  marginparwidth=0pt,
  marginparsep=0pt
]{geometry}

% ─── LANGUAGE & ENCODING ─────────────────────────────────────────────────────
\usepackage[utf8]{inputenc}              % UTF-8 input
\usepackage[T1]{fontenc}                 % T1 font encoding
\usepackage[english]{babel}              % English hyphenation
% Add arabic option if bilingual content: \usepackage[english,arabic]{babel}

% ─── FONTS ───────────────────────────────────────────────────────────────────
% Times New Roman: universal thesis standard (Oxford, Cambridge, APA, Chicago,
% MIT, Stanford, FCI Cairo, Ain Shams)
\usepackage{fontspec}                    % XeLaTeX / LuaLaTeX font selection
\setmainfont{Times New Roman}            % Body serif font — 12pt standard
\setsansfont{Arial}                      % Sans-serif for captions, headers
\setmonofont{Courier New}[Scale=0.90]    % Monospace for code — 10~11pt equiv.

% Fallback if TNR not available:
% \setmainfont{TeX Gyre Termes}         % Free Times clone
% \setsansfont{TeX Gyre Heros}          % Free Helvetica/Arial clone
% \setmonofont{TeX Gyre Cursor}         % Free Courier clone

% ─── LINE SPACING ────────────────────────────────────────────────────────────
% FCI Cairo: 1.5 line spacing
% Oxford/Cambridge: 1.5 or double
% APA: double
% MIT/Stanford: 1.5 or double
% → 1.5 is the most common for FCI; use \doublespacing if your dept. requires
\usepackage{setspace}
\onehalfspacing

% ─── PARAGRAPH FORMATTING ────────────────────────────────────────────────────
% APA / Chicago / FCI: 0.5" (12.7mm) first-line indent
% APA: NO extra space between paragraphs
\usepackage{indentfirst}                 % indent first paragraph after heading
\setlength{\parindent}{12.7mm}           % 0.5 inch first-line indent
\setlength{\parskip}{0pt}                % no extra vertical space between paras

% ─── SECTIONAL HEADINGS ─────────────────────────────────────────────────────
% Heading hierarchy based on APA 7th + FCI Cairo synthesis:
%   Chapter  → centered, bold, 16pt (FCI standard)
%   Section  → left-aligned, bold, 14pt (FCI standard)
%   Subsection → left-aligned, bold, 12pt (FCI standard)
%   Subsubsection → left-aligned, bold italic, 12pt
%   Paragraph → indented, bold, run-in, period, 12pt
\usepackage{titlesec}
\usepackage{titletoc}

% ---- Chapter (APA L1 equivalent) ----
% FCI: bold, centered, 16pt
% Oxford/Cambridge: centered, bold, 14-16pt
\titleformat{\chapter}[display]
  {\normalfont\centering\bfseries}
  {\MakeUppercase{\chaptertitlename}\ \thechapter}
  {20pt}
  {\fontsize{16pt}{19.2pt}\selectfont\bfseries}
\titlespacing*{\chapter}{0pt}{-10pt}{30pt}

% ---- Section (APA L2 equivalent) ----
% FCI: bold, left-aligned, 14pt
% APA: Flush Left, Bold, Title Case
\titleformat{\section}
  {\normalfont\bfseries\fontsize{14pt}{16.8pt}\selectfont}
  {\thesection}
  {1em}
  {}
\titlespacing*{\section}{0pt}{18pt plus 4pt minus 2pt}{10pt plus 2pt minus 2pt}

% ---- Subsection (APA L3 equivalent) ----
% FCI: bold, left-aligned, 12pt
% APA: Flush Left, Bold Italic, Title Case
\titleformat{\subsection}
  {\normalfont\bfseries\fontsize{12pt}{14.4pt}\selectfont}
  {\thesubsection}
  {1em}
  {}
\titlespacing*{\subsection}{0pt}{14pt plus 4pt minus 2pt}{8pt plus 2pt minus 2pt}

% ---- Subsubsection (APA L4 equivalent) ----
% APA: Indented, Bold, Title Case, Period., Text continues
% Chicago: Flush Left, Italic
\titleformat{\subsubsection}
  {\normalfont\bfseries\itshape}
  {\thesubsubsection}
  {1em}
  {}
\titlespacing*{\subsubsection}{0pt}{12pt plus 4pt minus 2pt}{6pt plus 2pt minus 2pt}

% ---- Paragraph (APA L5 equivalent) ----
% APA: Indented, Bold Italic, Title Case, Period., Text continues
\titleformat{\paragraph}[runin]
  {\normalfont\bfseries\itshape}
  {\theparagraph}
  {1em}
  {}
\titlespacing*{\paragraph}{\parindent}{10pt plus 4pt minus 2pt}{1em plus 2pt minus 2pt}

% ─── TABLE OF CONTENTS FORMATTING ───────────────────────────────────────────
% Oxford-style: clean, professional dot leaders
\titlecontents{chapter}
  [0em]
  {\bfseries\vspace{6pt}}
  {\contentslabel{2em}}
  {\hspace*{2em}}
  {\titlerule*[8pt]{$\cdot$}\contentspage}

\titlecontents{section}
  [2em]
  {}
  {\contentslabel{2.5em}}
  {\hspace*{2.5em}}
  {\titlerule*[8pt]{$\cdot$}\contentspage}

\titlecontents{subsection}
  [4.5em]
  {}
  {\contentslabel{3em}}
  {\hspace*{3em}}
  {\titlerule*[8pt]{$\cdot$}\contentspage}

\titlecontents{subsubsection}
  [7.5em]
  {\small}
  {\contentslabel{3.5em}}
  {\hspace*{3.5em}}
  {\titlerule*[8pt]{$\cdot$}\contentspage}

% ─── NUMBERING DEPTH ────────────────────────────────────────────────────────
\setcounter{secnumdepth}{4}              % number up to subsubsection
\setcounter{tocdepth}{3}                 % show up to subsection in TOC
                                          % (change to 2 for sections only)

% ─── FIGURE & TABLE NUMBERING ───────────────────────────────────────────────
% Per-chapter numbering: Fig 3.1, Table 3.2 (MIT, Oxford, FCI standard)
\usepackage{chngcntr}
\counterwithin{figure}{chapter}
\counterwithin{table}{chapter}
% For continuous numbering (Fig 1, Fig 2, ...), use instead:
% \counterwithout{figure}{chapter}
% \counterwithout{table}{chapter}

% ─── PACKAGES — TABLES ──────────────────────────────────────────────────────
% APA 7th: ONLY 3 horizontal rules — top, header-body separator, bottom
% NO vertical rules. This is the universal academic standard.
\usepackage{booktabs}                    % \toprule, \midrule, \bottomrule
\usepackage{tabularx}                    % flexible-width tables
\usepackage{longtable}                   % tables spanning multiple pages
\usepackage{multirow}                    % merged cells vertically
\usepackage{makecell}                    % line breaks inside cells
\usepackage{array}                       % extended column types
\usepackage{colortbl}                    % colored cells (use sparingly)

% ─── PACKAGES — FIGURES ─────────────────────────────────────────────────────
\usepackage{graphicx}                    % \includegraphics
\usepackage{float}                       % [H] placement specifier
\usepackage{subcaption}                  % sub-figures (a), (b), (c)
\usepackage{wrapfig}                     % text wrapping around figures
\usepackage{tikz}
\usetikzlibrary{
  shapes, arrows, positioning, calc,
  fit, backgrounds, decorations.pathreplacing,
  matrix, shadows, patterns
}

% ─── PACKAGES — CAPTIONS ────────────────────────────────────────────────────
% Universal standard: Table caption ABOVE, Figure caption BELOW
% APA: Table title italic, "Note." below table
\usepackage[
  font={small},
  labelfont={bf},
  labelsep=quad,
  justification=centerlast,
  skip=8pt
]{caption}

% Figure caption: below figure (ALL standards agree)
\captionsetup[figure]{
  font={small},
  labelfont={bf},
  labelsep=quad,
  justification=centerlast,
  skip=6pt,
  position=bottom
}

% Table caption: above table (ALL standards agree)
\captionsetup[table]{
  font={small},
  labelfont={bf},
  labelsep=quad,
  justification=centerlast,
  skip=6pt,
  position=top
}

% Sub-figure caption
\captionsetup[subfigure]{
  font={footnote},
  labelfont={bf},
  labelsep=space,
  justification=centerlast,
  skip=4pt
}

% ─── PACKAGES — HEADERS & FOOTERS ───────────────────────────────────────────
\usepackage{fancyhdr}

% ─── PACKAGES — COLORS ──────────────────────────────────────────────────────
\usepackage{xcolor}
\usepackage[table]{xcolor}

% ─── PACKAGES — MATHEMATICS ─────────────────────────────────────────────────
\usepackage{amsmath}
\usepackage{amssymb}
\usepackage{amsfonts}
\usepackage{mathtools}
\usepackage{bm}

% ─── PACKAGES — ALGORITHMS ──────────────────────────────────────────────────
\usepackage{algorithm}
\usepackage{algpseudocode}
\usepackage{algorithmic}

% ─── PACKAGES — CODE LISTINGS ───────────────────────────────────────────────
\usepackage{listings}

% ─── PACKAGES — HYPERLINKS ──────────────────────────────────────────────────
\usepackage[
  colorlinks=true,
  linkcolor=ThesisBlue,
  citecolor=ThesisGreen,
  urlcolor=ThesisNavy,
  bookmarks=true,
  bookmarksopen=true,
  bookmarksnumbered=true,
  pdfauthor={\AuthorName},
  pdftitle={\ThesisTitle}
]{hyperref}

% ─── PACKAGES — BIBLIOGRAPHY ────────────────────────────────────────────────
% IEEE style: standard for CS departments (FCI, Ain Shams, MIT EECS)
\usepackage[backend=biber, style=ieee]{biblatex}
% \addbibresource{references.bib}        % uncomment and set your .bib file

% ─── PACKAGES — GLOSSARY / ACRONYMS ─────────────────────────────────────────
\usepackage[acronym, toc]{glossaries}
% \makeglossaries

% ─── PACKAGES — MISCELLANEOUS ───────────────────────────────────────────────
\usepackage{enumitem}                    % customized lists
\usepackage{pifont}                      % ding symbols
\usepackage{calc}
\usepackage{ifthen}
\usepackage{afterpage}
\usepackage{emptypage}                   % blank pages have no headers/footers
\usepackage{appendix}
\usepackage{pdflscape}                   % landscape pages
\usepackage{placeins}                    % \FloatBarrier
\usepackage{siunitx}                     % SI units
\usepackage{tcolorbox}                   % colored boxes for callouts
\usepackage{csquotes}                    % proper quoting (required by biblatex)

% ─── COLOR PALETTE ──────────────────────────────────────────────────────────
% Muted academic colors — professional, not distracting
% Inspired by Oxford/Cambridge thesis templates

% ---- Primary ----
\definecolor{ThesisBlue}{RGB}{0, 70, 140}         % headings, links
\definecolor{ThesisNavy}{RGB}{0, 40, 85}           % deep emphasis
\definecolor{ThesisGreen}{RGB}{0, 120, 60}         % citations, positive
\definecolor{ThesisRed}{RGB}{180, 30, 30}           % warnings, negative
\definecolor{ThesisOrange}{RGB}{210, 105, 0}        % highlights

% ---- Table ----
\definecolor{TableHeader}{RGB}{0, 70, 140}          % header row background
\definecolor{TableHeaderText}{RGB}{255, 255, 255}   % header row text
\definecolor{TableRowA}{RGB}{240, 245, 250}         % alternating row A
\definecolor{TableRowB}{RGB}{255, 255, 255}         % alternating row B

% ---- Code ----
\definecolor{CodeBG}{RGB}{248, 248, 248}
\definecolor{CodeFrame}{RGB}{200, 200, 200}
\definecolor{CodeKeyword}{RGB}{0, 0, 180}
\definecolor{CodeString}{RGB}{0, 128, 0}
\definecolor{CodeComment}{RGB}{128, 128, 128}
\definecolor{CodeNumber}{RGB}{160, 100, 0}

% ---- Callout Boxes ----
\definecolor{NoteBoxBG}{RGB}{232, 244, 255}
\definecolor{NoteBoxFrame}{RGB}{0, 70, 140}
\definecolor{WarningBoxBG}{RGB}{255, 243, 224}
\definecolor{WarningBoxFrame}{RGB}{210, 105, 0}
\definecolor{ImportantBoxBG}{RGB}{255, 232, 232}
\definecolor{ImportantBoxFrame}{RGB}{180, 30, 30}

% ─── PAGE STYLE — HEADERS & FOOTERS ────────────────────────────────────────
% Oxford/Cambridge: chapter title left header, page number right
% APA: page number top right (student paper) or running head + page# (pro)
% FCI: page number bottom center is common
% → Using Oxford-style with page number in header

\pagestyle{fancy}
\fancyhf{}

% ---- Header ----
\fancyhead[LE,RO]{\small\thepage}                     % page number
\fancyhead[RE]{\small\itshape\nouppercase{\leftmark}}  % chapter title (even)
\fancyhead[LO]{}                                       % leave blank (odd)
\renewcommand{\headrulewidth}{0.4pt}

% ---- Footer ----
\fancyfoot[C]{}
\fancyfoot[L]{\tiny\textcolor{gray}{\UniversityNameEN}}

% ---- Chapter-opening pages: no header, page number bottom center ----
\fancypagestyle{plain}{%
  \fancyhf{}%
  \fancyfoot[C]{\small\thepage}%
  \renewcommand{\headrulewidth}{0pt}%
}

% ─── TABLE FORMATTING ───────────────────────────────────────────────────────
% APA 7th (universally accepted):
%   - ONLY 3 horizontal rules: \toprule, \midrule, \bottomrule
%   - NO vertical rules
%   - Caption ABOVE table
%   - "Note." line below table (optional)
% Table text may be smaller (10pt) than body (12pt) — APA, Chicago

\renewcommand{\arraystretch}{1.30}       % row height (APA recommends 1.25-1.35)
\tabcolsep=8pt                           % column separation

% ---- Custom column types ----
\newcolumntype{L}[1]{>{\raggedright\arraybackslash}p{#1}}
\newcolumntype{C}[1]{>{\centering\arraybackslash}p{#1}}
\newcolumntype{R}[1]{>{\raggedleft\arraybackslash}p{#1}}
\newcolumntype{M}[1]{>{\centering\arraybackslash}m{#1}}

% ---- Styled table header command ----
% Use within booktabs tables for professional APA-style headers
\newcommand{\thead}[1]{\textbf{#1}}

% ---- Table Note command (APA style) ----
% Usage: \tablenote{Source: Author's calculations. n = 200.}
\newcommand{\tablenote}[1]{%
  \vspace{2pt}%
  {\small\textit{Note.}\quad #1}%
}

% ---- Professional table environment ----
% Usage:
%   \begin{proftable}{caption}{label}{col spec}
%     \toprule
%     \thead{Col 1} & \thead{Col 2} \\
%     \midrule
%     data & data \\
%     \bottomrule
%   \end{proftable}
\newenvironment{proftable}[3]{%
  \begin{table}[H]%
  \centering%
  \caption{#1}%
  \label{#2}%
  \small%
  \begin{tabular}{#3}%
}{%
  \end{tabular}%
  \end{table}%
}

% ---- Long table environment (multi-page) ----
% Same rules apply: booktabs only, no vertical rules
\newenvironment{proflongtable}[3]{%
  \begin{longtable}{#3}%
  \caption{#1}\label{#2}\\%
}{%
  \end{longtable}%
}

% ─── FIGURE PLACEMENT ───────────────────────────────────────────────────────
% APA: figures embedded near first mention, not floating far away
\renewcommand{\topfraction}{0.92}
\renewcommand{\bottomfraction}{0.85}
\renewcommand{\textfraction}{0.08}
\renewcommand{\floatpagefraction}{0.90}

% ─── LIST FORMATTING ────────────────────────────────────────────────────────
% APA: bullet lists allowed, no extra spacing between items
% Chicago: consistent indentation

\setitemize{
  leftmargin=2em,
  itemsep=4pt,
  parsep=2pt,
  topsep=6pt
}

\setenumerate{
  leftmargin=2em,
  itemsep=4pt,
  parsep=2pt,
  topsep=6pt,
  label=\arabic*.
}

\setdescription{
  leftmargin=2em,
  itemsep=4pt,
  parsep=2pt,
  topsep=6pt,
  font=\normalfont\bfseries
}

% ─── CODE LISTING FORMATTING ────────────────────────────────────────────────
% FCI Cairo: Courier New 10pt for code
\lstdefinestyle{ThesisCode}{
  backgroundcolor=\color{CodeBG},
  basicstyle=\ttfamily\footnotesize,
  keywordstyle=\color{CodeKeyword}\bfseries,
  stringstyle=\color{CodeString},
  commentstyle=\color{CodeComment}\itshape,
  numberstyle=\tiny\color{CodeNumber},
  numbers=left,
  numbersep=8pt,
  frame=single,
  rulecolor=\color{CodeFrame},
  framesep=5pt,
  xleftmargin=15pt,
  xrightmargin=5pt,
  breaklines=true,
  showstringspaces=false,
  tabsize=4,
  captionpos=b,
  columns=flexible,
  keepspaces=true,
  escapeinside={(*@}{@*)}
}
\lstset{style=ThesisCode}

\lstdefinestyle{PythonCode}{style=ThesisCode, language=Python}
\lstdefinestyle{BashCode}{style=ThesisCode, language=bash}

\newcommand{\code}[1]{\texttt{\small #1}}

% ─── ALGORITHM / PSEUDOCODE FORMATTING ──────────────────────────────────────
\algrenewcommand\algorithmicrequire{\textbf{Input:}}
\algrenewcommand\algorithmicensure{\textbf{Output:}}
\floatname{algorithm}{Algorithm}
\renewcommand{\algorithmiccomment}[1]{\hfill$\triangleright$ \textit{#1}}

% ─── MATHEMATICS FORMATTING ─────────────────────────────────────────────────
% Per-chapter equation numbering: (3.1), (3.2) — MIT, Oxford, Springer standard
\DeclareMathOperator*{\argmin}{arg\,min}
\DeclareMathOperator*{\argmax}{arg\,max}
\DeclareMathOperator{\KL}{KL}
\DeclareMathOperator{\softmax}{softmax}
\DeclareMathOperator{\sigmoid}{sigmoid}
\DeclareMathOperator{\ReLU}{ReLU}

% Math shortcuts
\newcommand{\R}{\mathbb{R}}
\newcommand{\E}{\mathbb{E}}
\newcommand{\norm}[1]{\left\lVert #1 \right\rVert}
\newcommand{\abs}[1]{\left\lvert #1 \right\rvert}
\newcommand{\inner}[2]{\left\langle #1,\, #2 \right\rangle}
\newcommand{\condbar}{\,\vert\,}
\newcommand{\bigO}{\mathcal{O}}

\numberwithin{equation}{chapter}

% ─── CALLOUT BOXES ──────────────────────────────────────────────────────────
% Use sparingly — academic theses prefer inline text over boxes
% But useful for: definitions, key findings, warnings

\newtcolorbox{notebox}[1][Note]{%
  colback=NoteBoxBG,
  colframe=NoteBoxFrame,
  fonttitle=\bfseries\sffamily,
  title=#1,
  boxrule=0.8pt,
  arc=2pt,
  left=8pt, right=8pt, top=4pt, bottom=4pt,
  before skip=10pt, after skip=10pt
}

\newtcolorbox{warningbox}[1][Warning]{%
  colback=WarningBoxBG,
  colframe=WarningBoxFrame,
  fonttitle=\bfseries\sffamily,
  title=#1,
  boxrule=0.8pt,
  arc=2pt,
  left=8pt, right=8pt, top=4pt, bottom=4pt,
  before skip=10pt, after skip=10pt
}

\newtcolorbox{importantbox}[1][Important]{%
  colback=ImportantBoxBG,
  colframe=ImportantBoxFrame,
  fonttitle=\bfseries\sffamily,
  title=#1,
  boxrule=0.8pt,
  arc=2pt,
  left=8pt, right=8pt, top=4pt, bottom=4pt,
  before skip=10pt, after skip=10pt
}

\newtcolorbox{definitionbox}[1][Definition]{%
  colback=NoteBoxBG,
  colframe=ThesisBlue,
  fonttitle=\bfseries\sffamily\color{white},
  coltitle=white,
  colbacktitle=ThesisBlue,
  title=#1,
  boxrule=0.8pt,
  arc=2pt,
  left=8pt, right=8pt, top=4pt, bottom=4pt,
  before skip=10pt, after skip=10pt
}

% ─── CROSS-REFERENCE COMMANDS ───────────────────────────────────────────────
% Friendly references — universal academic style
\newcommand{\reffig}[1]{Figure~\ref{#1}}
\newcommand{\reftab}[1]{Table~\ref{#1}}
\newcommand{\refeq}[1]{Equation~\ref{#1}}
\newcommand{\refch}[1]{Chapter~\ref{#1}}
\newcommand{\refsec}[1]{Section~\ref{#1}}
\newcommand{\refalg}[1]{Algorithm~\ref{#1}}
\newcommand{\refapp}[1]{Appendix~\ref{#1}}
\newcommand{\reflst}[1]{Listing~\ref{#1}}
\newcommand{\refpage}[1]{page~\pageref{#1}}

% ─── UTILITY COMMANDS ───────────────────────────────────────────────────────
\newcommand{\pct}[1]{#1\%}
\newcommand{\metric}[2]{\textbf{#1}\,=\,#2}

% ─── COVER PAGE ─────────────────────────────────────────────────────────────
% FCI Cairo style: University name, Faculty, Department, Title, Author, Supervisor
\newcommand{\MakeCoverPage}{%
  \begin{titlepage}
    \centering
    \vspace*{0.5cm}

    {\scshape\LARGE \UniversityFullEN \par}
    \vspace{0.3cm}
    {\scshape\Large \UniversityNameEN \par}
    \vspace{0.2cm}
    {\scshape\large \DepartmentNameEN \par}

    \vspace{2cm}

    {\fontsize{18pt}{21.6pt}\selectfont\bfseries \ThesisTitle \par}

    \vspace{1.5cm}

    {\Large Graduation Project Thesis\par}

    \vspace{2cm}

    \begin{minipage}{0.45\textwidth}
      \begin{flushleft}\large
        \textbf{Submitted by:}\\
        \AuthorName
      \end{flushleft}
    \end{minipage}
    \hfill
    \begin{minipage}{0.45\textwidth}
      \begin{flushright}\large
        \textbf{Supervised by:}\\
        \SupervisorName
      \end{flushright}
    \end{minipage}

    \vfill

    {\large \AcademicYear\par}

  \end{titlepage}
}

% ─── ABSTRACT PAGE ──────────────────────────────────────────────────────────
% FCI: Arabic abstract first, then English (or vice versa per department)
\newcommand{\MakeAbstractPage}[2]{%
  \chapter*{Abstract}
  \addcontentsline{toc}{chapter}{Abstract}
  #1
  \vfill
  \ifthenelse{\equal{#2}{}}{}{%
    \newpage
    \chapter*{الملخص}
    \addcontentsline{toc}{chapter}{الملخص}
    #2
  }
}

% ─── ACKNOWLEDGMENTS ────────────────────────────────────────────────────────
\newcommand{\MakeAcknowledgmentsPage}[1]{%
  \chapter*{Acknowledgments}
  \addcontentsline{toc}{chapter}{Acknowledgments}
  #1
}

% ─── DEDICATION ─────────────────────────────────────────────────────────────
\newcommand{\MakeDedicationPage}[1]{%
  \vspace*{\fill}
  \begin{center}
    \textit{\Large #1}
  \end{center}
  \vspace*{\fill}
}

% ─── APPENDIX FORMATTING ────────────────────────────────────────────────────
% APA: labeled Appendix A, Appendix B; each on new page
% FCI: similar structure
\renewcommand{\appendixname}{Appendix}

% ─── BIBLIOGRAPHY FORMATTING ────────────────────────────────────────────────
% IEEE style for CS departments (FCI, Ain Shams, MIT EECS)
\defbibheading{bibliography}{%
  \chapter*{Bibliography}%
  \addcontentsline{toc}{chapter}{Bibliography}%
}

% ─── FRONT MATTER / MAIN MATTER ─────────────────────────────────────────────
% MIT/Oxford/Cambridge standard:
%   Front matter: Roman numerals (i, ii, iii, ...)
%   Main matter:  Arabic numerals (1, 2, 3, ...)
%   Chapter 1 starts at page 1
\makeatletter
\renewcommand\frontmatter{%
  \cleardoublepage
  \pagenumbering{roman}
}

\renewcommand\mainmatter{%
  \cleardoublepage
  \pagenumbering{arabic}
}

\renewcommand\backmatter{%
  \cleardoublepage
  \setcounter{chapter}{0}%
}
\makeatother

% ─── DRAFT MARKERS ──────────────────────────────────────────────────────────
\newif\ifdraft
\drafttrue                              % set to \draftfalse for final version

\ifdraft
  \newcommand{\todo}[1]{\textcolor{ThesisRed}{\bfseries[TODO: #1]}}
  \newcommand{\FIXME}[1]{\textcolor{ThesisRed}{\bfseries[FIXME: #1]}}
  \newcommand{\CITE}[1]{\textcolor{ThesisOrange}{\bfseries[CITE: #1]}}
\else
  \newcommand{\todo}[1]{}
  \newcommand{\FIXME}[1]{}
  \newcommand{\CITE}[1]{\cite{#1}}
\fi

% ─── ORPHAN / WIDOW PREVENTION ──────────────────────────────────────────────
% Chicago/Oxford: avoid orphan/widow lines
\clubpenalty=10000
\widowpenalty=10000
\emergencystretch=1em
\tolerance=200
\setlength{\emergencystretch}{3em}

% ─── HYPHENATION EXCEPTIONS ─────────────────────────────────────────────────
\hyphenation{
  fine-tun-ing
  know-ledge
  distil-la-tion
  eval-ua-tion
  pre-train-ed
  multi-ple
  parameter
  hyper-parameter
  archi-tec-ture
}

% ─── SI UNIT FORMATTING ─────────────────────────────────────────────────────
\sisetup{
  detect-all,
  group-separator={,},
  group-minimum-digits=4,
  per-mode=symbol
}

% =============================================================================
% FORMATTING DECISIONS LOG — Based on Cross-Source Research
% =============================================================================
%
% | Decision              | Chosen Value       | Sources Supporting It          |
% |-----------------------|--------------------|--------------------------------|
% | Page size             | A4                 | Oxford, Cambridge, Cairo, FCI  |
% | Inner margin          | 3.5cm              | FCI (3cm), MIT (1.5"), avg     |
% | Outer margin          | 2.5cm              | FCI, Egyptian Council          |
% | Body font             | TNR 12pt           | ALL sources agree              |
% | Line spacing          | 1.5                | FCI, Cambridge, Oxford         |
% | Paragraph indent      | 12.7mm (0.5")      | APA, Chicago, FCI              |
% | Chapter heading       | 16pt bold centered | FCI, Oxford, Cambridge         |
% | Section heading       | 14pt bold left     | FCI standard                   |
% | Subsection heading    | 12pt bold left     | FCI, APA L3                    |
% | Table rules           | 3 only (booktabs)  | APA 7th, ALL academic stds     |
% | Table caption         | Above table        | APA, IEEE, ACM, ALL sources    |
% | Figure caption        | Below figure       | APA, IEEE, ACM, ALL sources    |
% | Equation numbering    | Per chapter        | MIT, Oxford, Springer          |
% | Figure numbering      | Per chapter        | MIT, Oxford, FCI               |
% | Bibliography style    | IEEE               | FCI CS dept, ACM, MIT EECS     |
% | Front matter pages    | Roman numerals     | MIT, Oxford, Cambridge, FCI    |
% | Main matter pages     | Arabic numerals    | ALL sources agree              |
% | TOC depth             | 3 levels           | Oxford, Cambridge              |
% | Secnum depth          | 4 levels           | Standard for theses            |
%
% =============================================================================

% =============================================================================
% END OF FORMATTING FILE v2
% =============================================================================
 in main.tex preamble
% =============================================================================

% ─── DOCUMENT CLASS ──────────────────────────────────────────────────────────
% book class: best for multi-chapter theses (front/main/back matter support)
% 12pt: universal standard (Oxford, Cambridge, APA, Chicago, Cairo Univ)
% a4paper: international standard (Oxford, Cambridge, Egyptian universities)
% oneside: standard for single-sided printing
% openright: chapters start on right-hand (odd) pages — MIT/Oxford standard
\documentclass[12pt, a4paper, oneside, openright]{report}
% For two-sided printing (required by some universities), change to:
% \documentclass[12pt, a4paper, twoside, openright]{report}

% ─── USER CONFIGURATION ─────────────────────────────────────────────────────
\newcommand{\UniversityNameEN}{Faculty of Computers and Information}
\newcommand{\UniversityNameAR}{كلية الحاسبات والمعلومات}
\newcommand{\DepartmentNameEN}{Computer Science Department}
\newcommand{\DepartmentNameAR}{قسم علوم الحاسب}
\newcommand{\UniversityFullEN}{Cairo University}
\newcommand{\ThesisTitle}{Quiz Generation via Knowledge Distillation and Fine-Tuning}
\newcommand{\AuthorName}{Rehab Hamdy}
\newcommand{\SupervisorName}{Prof.\ XXX}
\newcommand{\AcademicYear}{2025--2026}

% ─── PAGE GEOMETRY ───────────────────────────────────────────────────────────
% FCI Cairo: 3cm left (binding), 2.5cm others
% Egyptian Supreme Council: 4cm binding side, 2.5cm others
% MIT/Stanford: 1.5" (38mm) inner for binding
% Compromise: 3.5cm inner (binding), 2.5cm others — satisfies all standards
\usepackage[
  a4paper,
  top=2.5cm,
  bottom=2.5cm,
  inner=3.5cm,                          % binding side (left for English)
  outer=2.5cm,
  headheight=14.5pt,
  headsep=15mm,
  footskip=20mm,
  marginparwidth=0pt,
  marginparsep=0pt
]{geometry}

% ─── LANGUAGE & ENCODING ─────────────────────────────────────────────────────
\usepackage[utf8]{inputenc}              % UTF-8 input
\usepackage[T1]{fontenc}                 % T1 font encoding
\usepackage[english]{babel}              % English hyphenation
% Add arabic option if bilingual content: \usepackage[english,arabic]{babel}

% ─── FONTS ───────────────────────────────────────────────────────────────────
% Times New Roman: universal thesis standard (Oxford, Cambridge, APA, Chicago,
% MIT, Stanford, FCI Cairo, Ain Shams)
\usepackage{fontspec}                    % XeLaTeX / LuaLaTeX font selection
\setmainfont{Times New Roman}            % Body serif font — 12pt standard
\setsansfont{Arial}                      % Sans-serif for captions, headers
\setmonofont{Courier New}[Scale=0.90]    % Monospace for code — 10~11pt equiv.

% Fallback if TNR not available:
% \setmainfont{TeX Gyre Termes}         % Free Times clone
% \setsansfont{TeX Gyre Heros}          % Free Helvetica/Arial clone
% \setmonofont{TeX Gyre Cursor}         % Free Courier clone

% ─── LINE SPACING ────────────────────────────────────────────────────────────
% FCI Cairo: 1.5 line spacing
% Oxford/Cambridge: 1.5 or double
% APA: double
% MIT/Stanford: 1.5 or double
% → 1.5 is the most common for FCI; use \doublespacing if your dept. requires
\usepackage{setspace}
\onehalfspacing

% ─── PARAGRAPH FORMATTING ────────────────────────────────────────────────────
% APA / Chicago / FCI: 0.5" (12.7mm) first-line indent
% APA: NO extra space between paragraphs
\usepackage{indentfirst}                 % indent first paragraph after heading
\setlength{\parindent}{12.7mm}           % 0.5 inch first-line indent
\setlength{\parskip}{0pt}                % no extra vertical space between paras

% ─── SECTIONAL HEADINGS ─────────────────────────────────────────────────────
% Heading hierarchy based on APA 7th + FCI Cairo synthesis:
%   Chapter  → centered, bold, 16pt (FCI standard)
%   Section  → left-aligned, bold, 14pt (FCI standard)
%   Subsection → left-aligned, bold, 12pt (FCI standard)
%   Subsubsection → left-aligned, bold italic, 12pt
%   Paragraph → indented, bold, run-in, period, 12pt
\usepackage{titlesec}
\usepackage{titletoc}

% ---- Chapter (APA L1 equivalent) ----
% FCI: bold, centered, 16pt
% Oxford/Cambridge: centered, bold, 14-16pt
\titleformat{\chapter}[display]
  {\normalfont\centering\bfseries}
  {\MakeUppercase{\chaptertitlename}\ \thechapter}
  {20pt}
  {\fontsize{16pt}{19.2pt}\selectfont\bfseries}
\titlespacing*{\chapter}{0pt}{-10pt}{30pt}

% ---- Section (APA L2 equivalent) ----
% FCI: bold, left-aligned, 14pt
% APA: Flush Left, Bold, Title Case
\titleformat{\section}
  {\normalfont\bfseries\fontsize{14pt}{16.8pt}\selectfont}
  {\thesection}
  {1em}
  {}
\titlespacing*{\section}{0pt}{18pt plus 4pt minus 2pt}{10pt plus 2pt minus 2pt}

% ---- Subsection (APA L3 equivalent) ----
% FCI: bold, left-aligned, 12pt
% APA: Flush Left, Bold Italic, Title Case
\titleformat{\subsection}
  {\normalfont\bfseries\fontsize{12pt}{14.4pt}\selectfont}
  {\thesubsection}
  {1em}
  {}
\titlespacing*{\subsection}{0pt}{14pt plus 4pt minus 2pt}{8pt plus 2pt minus 2pt}

% ---- Subsubsection (APA L4 equivalent) ----
% APA: Indented, Bold, Title Case, Period., Text continues
% Chicago: Flush Left, Italic
\titleformat{\subsubsection}
  {\normalfont\bfseries\itshape}
  {\thesubsubsection}
  {1em}
  {}
\titlespacing*{\subsubsection}{0pt}{12pt plus 4pt minus 2pt}{6pt plus 2pt minus 2pt}

% ---- Paragraph (APA L5 equivalent) ----
% APA: Indented, Bold Italic, Title Case, Period., Text continues
\titleformat{\paragraph}[runin]
  {\normalfont\bfseries\itshape}
  {\theparagraph}
  {1em}
  {}
\titlespacing*{\paragraph}{\parindent}{10pt plus 4pt minus 2pt}{1em plus 2pt minus 2pt}

% ─── TABLE OF CONTENTS FORMATTING ───────────────────────────────────────────
% Oxford-style: clean, professional dot leaders
\titlecontents{chapter}
  [0em]
  {\bfseries\vspace{6pt}}
  {\contentslabel{2em}}
  {\hspace*{2em}}
  {\titlerule*[8pt]{$\cdot$}\contentspage}

\titlecontents{section}
  [2em]
  {}
  {\contentslabel{2.5em}}
  {\hspace*{2.5em}}
  {\titlerule*[8pt]{$\cdot$}\contentspage}

\titlecontents{subsection}
  [4.5em]
  {}
  {\contentslabel{3em}}
  {\hspace*{3em}}
  {\titlerule*[8pt]{$\cdot$}\contentspage}

\titlecontents{subsubsection}
  [7.5em]
  {\small}
  {\contentslabel{3.5em}}
  {\hspace*{3.5em}}
  {\titlerule*[8pt]{$\cdot$}\contentspage}

% ─── NUMBERING DEPTH ────────────────────────────────────────────────────────
\setcounter{secnumdepth}{4}              % number up to subsubsection
\setcounter{tocdepth}{3}                 % show up to subsection in TOC
                                          % (change to 2 for sections only)

% ─── FIGURE & TABLE NUMBERING ───────────────────────────────────────────────
% Per-chapter numbering: Fig 3.1, Table 3.2 (MIT, Oxford, FCI standard)
\usepackage{chngcntr}
\counterwithin{figure}{chapter}
\counterwithin{table}{chapter}
% For continuous numbering (Fig 1, Fig 2, ...), use instead:
% \counterwithout{figure}{chapter}
% \counterwithout{table}{chapter}

% ─── PACKAGES — TABLES ──────────────────────────────────────────────────────
% APA 7th: ONLY 3 horizontal rules — top, header-body separator, bottom
% NO vertical rules. This is the universal academic standard.
\usepackage{booktabs}                    % \toprule, \midrule, \bottomrule
\usepackage{tabularx}                    % flexible-width tables
\usepackage{longtable}                   % tables spanning multiple pages
\usepackage{multirow}                    % merged cells vertically
\usepackage{makecell}                    % line breaks inside cells
\usepackage{array}                       % extended column types
\usepackage{colortbl}                    % colored cells (use sparingly)

% ─── PACKAGES — FIGURES ─────────────────────────────────────────────────────
\usepackage{graphicx}                    % \includegraphics
\usepackage{float}                       % [H] placement specifier
\usepackage{subcaption}                  % sub-figures (a), (b), (c)
\usepackage{wrapfig}                     % text wrapping around figures
\usepackage{tikz}
\usetikzlibrary{
  shapes, arrows, positioning, calc,
  fit, backgrounds, decorations.pathreplacing,
  matrix, shadows, patterns
}

% ─── PACKAGES — CAPTIONS ────────────────────────────────────────────────────
% Universal standard: Table caption ABOVE, Figure caption BELOW
% APA: Table title italic, "Note." below table
\usepackage[
  font={small},
  labelfont={bf},
  labelsep=quad,
  justification=centerlast,
  skip=8pt
]{caption}

% Figure caption: below figure (ALL standards agree)
\captionsetup[figure]{
  font={small},
  labelfont={bf},
  labelsep=quad,
  justification=centerlast,
  skip=6pt,
  position=bottom
}

% Table caption: above table (ALL standards agree)
\captionsetup[table]{
  font={small},
  labelfont={bf},
  labelsep=quad,
  justification=centerlast,
  skip=6pt,
  position=top
}

% Sub-figure caption
\captionsetup[subfigure]{
  font={footnote},
  labelfont={bf},
  labelsep=space,
  justification=centerlast,
  skip=4pt
}

% ─── PACKAGES — HEADERS & FOOTERS ───────────────────────────────────────────
\usepackage{fancyhdr}

% ─── PACKAGES — COLORS ──────────────────────────────────────────────────────
\usepackage{xcolor}
\usepackage[table]{xcolor}

% ─── PACKAGES — MATHEMATICS ─────────────────────────────────────────────────
\usepackage{amsmath}
\usepackage{amssymb}
\usepackage{amsfonts}
\usepackage{mathtools}
\usepackage{bm}

% ─── PACKAGES — ALGORITHMS ──────────────────────────────────────────────────
\usepackage{algorithm}
\usepackage{algpseudocode}
\usepackage{algorithmic}

% ─── PACKAGES — CODE LISTINGS ───────────────────────────────────────────────
\usepackage{listings}

% ─── PACKAGES — HYPERLINKS ──────────────────────────────────────────────────
\usepackage[
  colorlinks=true,
  linkcolor=ThesisBlue,
  citecolor=ThesisGreen,
  urlcolor=ThesisNavy,
  bookmarks=true,
  bookmarksopen=true,
  bookmarksnumbered=true,
  pdfauthor={\AuthorName},
  pdftitle={\ThesisTitle}
]{hyperref}

% ─── PACKAGES — BIBLIOGRAPHY ────────────────────────────────────────────────
% IEEE style: standard for CS departments (FCI, Ain Shams, MIT EECS)
\usepackage[backend=biber, style=ieee]{biblatex}
% \addbibresource{references.bib}        % uncomment and set your .bib file

% ─── PACKAGES — GLOSSARY / ACRONYMS ─────────────────────────────────────────
\usepackage[acronym, toc]{glossaries}
% \makeglossaries

% ─── PACKAGES — MISCELLANEOUS ───────────────────────────────────────────────
\usepackage{enumitem}                    % customized lists
\usepackage{pifont}                      % ding symbols
\usepackage{calc}
\usepackage{ifthen}
\usepackage{afterpage}
\usepackage{emptypage}                   % blank pages have no headers/footers
\usepackage{appendix}
\usepackage{pdflscape}                   % landscape pages
\usepackage{placeins}                    % \FloatBarrier
\usepackage{siunitx}                     % SI units
\usepackage{tcolorbox}                   % colored boxes for callouts
\usepackage{csquotes}                    % proper quoting (required by biblatex)

% ─── COLOR PALETTE ──────────────────────────────────────────────────────────
% Muted academic colors — professional, not distracting
% Inspired by Oxford/Cambridge thesis templates

% ---- Primary ----
\definecolor{ThesisBlue}{RGB}{0, 70, 140}         % headings, links
\definecolor{ThesisNavy}{RGB}{0, 40, 85}           % deep emphasis
\definecolor{ThesisGreen}{RGB}{0, 120, 60}         % citations, positive
\definecolor{ThesisRed}{RGB}{180, 30, 30}           % warnings, negative
\definecolor{ThesisOrange}{RGB}{210, 105, 0}        % highlights

% ---- Table ----
\definecolor{TableHeader}{RGB}{0, 70, 140}          % header row background
\definecolor{TableHeaderText}{RGB}{255, 255, 255}   % header row text
\definecolor{TableRowA}{RGB}{240, 245, 250}         % alternating row A
\definecolor{TableRowB}{RGB}{255, 255, 255}         % alternating row B

% ---- Code ----
\definecolor{CodeBG}{RGB}{248, 248, 248}
\definecolor{CodeFrame}{RGB}{200, 200, 200}
\definecolor{CodeKeyword}{RGB}{0, 0, 180}
\definecolor{CodeString}{RGB}{0, 128, 0}
\definecolor{CodeComment}{RGB}{128, 128, 128}
\definecolor{CodeNumber}{RGB}{160, 100, 0}

% ---- Callout Boxes ----
\definecolor{NoteBoxBG}{RGB}{232, 244, 255}
\definecolor{NoteBoxFrame}{RGB}{0, 70, 140}
\definecolor{WarningBoxBG}{RGB}{255, 243, 224}
\definecolor{WarningBoxFrame}{RGB}{210, 105, 0}
\definecolor{ImportantBoxBG}{RGB}{255, 232, 232}
\definecolor{ImportantBoxFrame}{RGB}{180, 30, 30}

% ─── PAGE STYLE — HEADERS & FOOTERS ────────────────────────────────────────
% Oxford/Cambridge: chapter title left header, page number right
% APA: page number top right (student paper) or running head + page# (pro)
% FCI: page number bottom center is common
% → Using Oxford-style with page number in header

\pagestyle{fancy}
\fancyhf{}

% ---- Header ----
\fancyhead[LE,RO]{\small\thepage}                     % page number
\fancyhead[RE]{\small\itshape\nouppercase{\leftmark}}  % chapter title (even)
\fancyhead[LO]{}                                       % leave blank (odd)
\renewcommand{\headrulewidth}{0.4pt}

% ---- Footer ----
\fancyfoot[C]{}
\fancyfoot[L]{\tiny\textcolor{gray}{\UniversityNameEN}}

% ---- Chapter-opening pages: no header, page number bottom center ----
\fancypagestyle{plain}{%
  \fancyhf{}%
  \fancyfoot[C]{\small\thepage}%
  \renewcommand{\headrulewidth}{0pt}%
}

% ─── TABLE FORMATTING ───────────────────────────────────────────────────────
% APA 7th (universally accepted):
%   - ONLY 3 horizontal rules: \toprule, \midrule, \bottomrule
%   - NO vertical rules
%   - Caption ABOVE table
%   - "Note." line below table (optional)
% Table text may be smaller (10pt) than body (12pt) — APA, Chicago

\renewcommand{\arraystretch}{1.30}       % row height (APA recommends 1.25-1.35)
\tabcolsep=8pt                           % column separation

% ---- Custom column types ----
\newcolumntype{L}[1]{>{\raggedright\arraybackslash}p{#1}}
\newcolumntype{C}[1]{>{\centering\arraybackslash}p{#1}}
\newcolumntype{R}[1]{>{\raggedleft\arraybackslash}p{#1}}
\newcolumntype{M}[1]{>{\centering\arraybackslash}m{#1}}

% ---- Styled table header command ----
% Use within booktabs tables for professional APA-style headers
\newcommand{\thead}[1]{\textbf{#1}}

% ---- Table Note command (APA style) ----
% Usage: \tablenote{Source: Author's calculations. n = 200.}
\newcommand{\tablenote}[1]{%
  \vspace{2pt}%
  {\small\textit{Note.}\quad #1}%
}

% ---- Professional table environment ----
% Usage:
%   \begin{proftable}{caption}{label}{col spec}
%     \toprule
%     \thead{Col 1} & \thead{Col 2} \\
%     \midrule
%     data & data \\
%     \bottomrule
%   \end{proftable}
\newenvironment{proftable}[3]{%
  \begin{table}[H]%
  \centering%
  \caption{#1}%
  \label{#2}%
  \small%
  \begin{tabular}{#3}%
}{%
  \end{tabular}%
  \end{table}%
}

% ---- Long table environment (multi-page) ----
% Same rules apply: booktabs only, no vertical rules
\newenvironment{proflongtable}[3]{%
  \begin{longtable}{#3}%
  \caption{#1}\label{#2}\\%
}{%
  \end{longtable}%
}

% ─── FIGURE PLACEMENT ───────────────────────────────────────────────────────
% APA: figures embedded near first mention, not floating far away
\renewcommand{\topfraction}{0.92}
\renewcommand{\bottomfraction}{0.85}
\renewcommand{\textfraction}{0.08}
\renewcommand{\floatpagefraction}{0.90}

% ─── LIST FORMATTING ────────────────────────────────────────────────────────
% APA: bullet lists allowed, no extra spacing between items
% Chicago: consistent indentation

\setitemize{
  leftmargin=2em,
  itemsep=4pt,
  parsep=2pt,
  topsep=6pt
}

\setenumerate{
  leftmargin=2em,
  itemsep=4pt,
  parsep=2pt,
  topsep=6pt,
  label=\arabic*.
}

\setdescription{
  leftmargin=2em,
  itemsep=4pt,
  parsep=2pt,
  topsep=6pt,
  font=\normalfont\bfseries
}

% ─── CODE LISTING FORMATTING ────────────────────────────────────────────────
% FCI Cairo: Courier New 10pt for code
\lstdefinestyle{ThesisCode}{
  backgroundcolor=\color{CodeBG},
  basicstyle=\ttfamily\footnotesize,
  keywordstyle=\color{CodeKeyword}\bfseries,
  stringstyle=\color{CodeString},
  commentstyle=\color{CodeComment}\itshape,
  numberstyle=\tiny\color{CodeNumber},
  numbers=left,
  numbersep=8pt,
  frame=single,
  rulecolor=\color{CodeFrame},
  framesep=5pt,
  xleftmargin=15pt,
  xrightmargin=5pt,
  breaklines=true,
  showstringspaces=false,
  tabsize=4,
  captionpos=b,
  columns=flexible,
  keepspaces=true,
  escapeinside={(*@}{@*)}
}
\lstset{style=ThesisCode}

\lstdefinestyle{PythonCode}{style=ThesisCode, language=Python}
\lstdefinestyle{BashCode}{style=ThesisCode, language=bash}

\newcommand{\code}[1]{\texttt{\small #1}}

% ─── ALGORITHM / PSEUDOCODE FORMATTING ──────────────────────────────────────
\algrenewcommand\algorithmicrequire{\textbf{Input:}}
\algrenewcommand\algorithmicensure{\textbf{Output:}}
\floatname{algorithm}{Algorithm}
\renewcommand{\algorithmiccomment}[1]{\hfill$\triangleright$ \textit{#1}}

% ─── MATHEMATICS FORMATTING ─────────────────────────────────────────────────
% Per-chapter equation numbering: (3.1), (3.2) — MIT, Oxford, Springer standard
\DeclareMathOperator*{\argmin}{arg\,min}
\DeclareMathOperator*{\argmax}{arg\,max}
\DeclareMathOperator{\KL}{KL}
\DeclareMathOperator{\softmax}{softmax}
\DeclareMathOperator{\sigmoid}{sigmoid}
\DeclareMathOperator{\ReLU}{ReLU}

% Math shortcuts
\newcommand{\R}{\mathbb{R}}
\newcommand{\E}{\mathbb{E}}
\newcommand{\norm}[1]{\left\lVert #1 \right\rVert}
\newcommand{\abs}[1]{\left\lvert #1 \right\rvert}
\newcommand{\inner}[2]{\left\langle #1,\, #2 \right\rangle}
\newcommand{\condbar}{\,\vert\,}
\newcommand{\bigO}{\mathcal{O}}

\numberwithin{equation}{chapter}

% ─── CALLOUT BOXES ──────────────────────────────────────────────────────────
% Use sparingly — academic theses prefer inline text over boxes
% But useful for: definitions, key findings, warnings

\newtcolorbox{notebox}[1][Note]{%
  colback=NoteBoxBG,
  colframe=NoteBoxFrame,
  fonttitle=\bfseries\sffamily,
  title=#1,
  boxrule=0.8pt,
  arc=2pt,
  left=8pt, right=8pt, top=4pt, bottom=4pt,
  before skip=10pt, after skip=10pt
}

\newtcolorbox{warningbox}[1][Warning]{%
  colback=WarningBoxBG,
  colframe=WarningBoxFrame,
  fonttitle=\bfseries\sffamily,
  title=#1,
  boxrule=0.8pt,
  arc=2pt,
  left=8pt, right=8pt, top=4pt, bottom=4pt,
  before skip=10pt, after skip=10pt
}

\newtcolorbox{importantbox}[1][Important]{%
  colback=ImportantBoxBG,
  colframe=ImportantBoxFrame,
  fonttitle=\bfseries\sffamily,
  title=#1,
  boxrule=0.8pt,
  arc=2pt,
  left=8pt, right=8pt, top=4pt, bottom=4pt,
  before skip=10pt, after skip=10pt
}

\newtcolorbox{definitionbox}[1][Definition]{%
  colback=NoteBoxBG,
  colframe=ThesisBlue,
  fonttitle=\bfseries\sffamily\color{white},
  coltitle=white,
  colbacktitle=ThesisBlue,
  title=#1,
  boxrule=0.8pt,
  arc=2pt,
  left=8pt, right=8pt, top=4pt, bottom=4pt,
  before skip=10pt, after skip=10pt
}

% ─── CROSS-REFERENCE COMMANDS ───────────────────────────────────────────────
% Friendly references — universal academic style
\newcommand{\reffig}[1]{Figure~\ref{#1}}
\newcommand{\reftab}[1]{Table~\ref{#1}}
\newcommand{\refeq}[1]{Equation~\ref{#1}}
\newcommand{\refch}[1]{Chapter~\ref{#1}}
\newcommand{\refsec}[1]{Section~\ref{#1}}
\newcommand{\refalg}[1]{Algorithm~\ref{#1}}
\newcommand{\refapp}[1]{Appendix~\ref{#1}}
\newcommand{\reflst}[1]{Listing~\ref{#1}}
\newcommand{\refpage}[1]{page~\pageref{#1}}

% ─── UTILITY COMMANDS ───────────────────────────────────────────────────────
\newcommand{\pct}[1]{#1\%}
\newcommand{\metric}[2]{\textbf{#1}\,=\,#2}

% ─── COVER PAGE ─────────────────────────────────────────────────────────────
% FCI Cairo style: University name, Faculty, Department, Title, Author, Supervisor
\newcommand{\MakeCoverPage}{%
  \begin{titlepage}
    \centering
    \vspace*{0.5cm}

    {\scshape\LARGE \UniversityFullEN \par}
    \vspace{0.3cm}
    {\scshape\Large \UniversityNameEN \par}
    \vspace{0.2cm}
    {\scshape\large \DepartmentNameEN \par}

    \vspace{2cm}

    {\fontsize{18pt}{21.6pt}\selectfont\bfseries \ThesisTitle \par}

    \vspace{1.5cm}

    {\Large Graduation Project Thesis\par}

    \vspace{2cm}

    \begin{minipage}{0.45\textwidth}
      \begin{flushleft}\large
        \textbf{Submitted by:}\\
        \AuthorName
      \end{flushleft}
    \end{minipage}
    \hfill
    \begin{minipage}{0.45\textwidth}
      \begin{flushright}\large
        \textbf{Supervised by:}\\
        \SupervisorName
      \end{flushright}
    \end{minipage}

    \vfill

    {\large \AcademicYear\par}

  \end{titlepage}
}

% ─── ABSTRACT PAGE ──────────────────────────────────────────────────────────
% FCI: Arabic abstract first, then English (or vice versa per department)
\newcommand{\MakeAbstractPage}[2]{%
  \chapter*{Abstract}
  \addcontentsline{toc}{chapter}{Abstract}
  #1
  \vfill
  \ifthenelse{\equal{#2}{}}{}{%
    \newpage
    \chapter*{الملخص}
    \addcontentsline{toc}{chapter}{الملخص}
    #2
  }
}

% ─── ACKNOWLEDGMENTS ────────────────────────────────────────────────────────
\newcommand{\MakeAcknowledgmentsPage}[1]{%
  \chapter*{Acknowledgments}
  \addcontentsline{toc}{chapter}{Acknowledgments}
  #1
}

% ─── DEDICATION ─────────────────────────────────────────────────────────────
\newcommand{\MakeDedicationPage}[1]{%
  \vspace*{\fill}
  \begin{center}
    \textit{\Large #1}
  \end{center}
  \vspace*{\fill}
}

% ─── APPENDIX FORMATTING ────────────────────────────────────────────────────
% APA: labeled Appendix A, Appendix B; each on new page
% FCI: similar structure
\renewcommand{\appendixname}{Appendix}

% ─── BIBLIOGRAPHY FORMATTING ────────────────────────────────────────────────
% IEEE style for CS departments (FCI, Ain Shams, MIT EECS)
\defbibheading{bibliography}{%
  \chapter*{Bibliography}%
  \addcontentsline{toc}{chapter}{Bibliography}%
}

% ─── FRONT MATTER / MAIN MATTER ─────────────────────────────────────────────
% MIT/Oxford/Cambridge standard:
%   Front matter: Roman numerals (i, ii, iii, ...)
%   Main matter:  Arabic numerals (1, 2, 3, ...)
%   Chapter 1 starts at page 1
\makeatletter
\renewcommand\frontmatter{%
  \cleardoublepage
  \pagenumbering{roman}
}

\renewcommand\mainmatter{%
  \cleardoublepage
  \pagenumbering{arabic}
}

\renewcommand\backmatter{%
  \cleardoublepage
  \setcounter{chapter}{0}%
}
\makeatother

% ─── DRAFT MARKERS ──────────────────────────────────────────────────────────
\newif\ifdraft
\drafttrue                              % set to \draftfalse for final version

\ifdraft
  \newcommand{\todo}[1]{\textcolor{ThesisRed}{\bfseries[TODO: #1]}}
  \newcommand{\FIXME}[1]{\textcolor{ThesisRed}{\bfseries[FIXME: #1]}}
  \newcommand{\CITE}[1]{\textcolor{ThesisOrange}{\bfseries[CITE: #1]}}
\else
  \newcommand{\todo}[1]{}
  \newcommand{\FIXME}[1]{}
  \newcommand{\CITE}[1]{\cite{#1}}
\fi

% ─── ORPHAN / WIDOW PREVENTION ──────────────────────────────────────────────
% Chicago/Oxford: avoid orphan/widow lines
\clubpenalty=10000
\widowpenalty=10000
\emergencystretch=1em
\tolerance=200
\setlength{\emergencystretch}{3em}

% ─── HYPHENATION EXCEPTIONS ─────────────────────────────────────────────────
\hyphenation{
  fine-tun-ing
  know-ledge
  distil-la-tion
  eval-ua-tion
  pre-train-ed
  multi-ple
  parameter
  hyper-parameter
  archi-tec-ture
}

% ─── SI UNIT FORMATTING ─────────────────────────────────────────────────────
\sisetup{
  detect-all,
  group-separator={,},
  group-minimum-digits=4,
  per-mode=symbol
}

% =============================================================================
% FORMATTING DECISIONS LOG — Based on Cross-Source Research
% =============================================================================
%
% | Decision              | Chosen Value       | Sources Supporting It          |
% |-----------------------|--------------------|--------------------------------|
% | Page size             | A4                 | Oxford, Cambridge, Cairo, FCI  |
% | Inner margin          | 3.5cm              | FCI (3cm), MIT (1.5"), avg     |
% | Outer margin          | 2.5cm              | FCI, Egyptian Council          |
% | Body font             | TNR 12pt           | ALL sources agree              |
% | Line spacing          | 1.5                | FCI, Cambridge, Oxford         |
% | Paragraph indent      | 12.7mm (0.5")      | APA, Chicago, FCI              |
% | Chapter heading       | 16pt bold centered | FCI, Oxford, Cambridge         |
% | Section heading       | 14pt bold left     | FCI standard                   |
% | Subsection heading    | 12pt bold left     | FCI, APA L3                    |
% | Table rules           | 3 only (booktabs)  | APA 7th, ALL academic stds     |
% | Table caption         | Above table        | APA, IEEE, ACM, ALL sources    |
% | Figure caption        | Below figure       | APA, IEEE, ACM, ALL sources    |
% | Equation numbering    | Per chapter        | MIT, Oxford, Springer          |
% | Figure numbering      | Per chapter        | MIT, Oxford, FCI               |
% | Bibliography style    | IEEE               | FCI CS dept, ACM, MIT EECS     |
% | Front matter pages    | Roman numerals     | MIT, Oxford, Cambridge, FCI    |
% | Main matter pages     | Arabic numerals    | ALL sources agree              |
% | TOC depth             | 3 levels           | Oxford, Cambridge              |
% | Secnum depth          | 4 levels           | Standard for theses            |
%
% =============================================================================

% =============================================================================
% END OF FORMATTING FILE v2
% =============================================================================
 in main.tex preamble
% =============================================================================

% ─── DOCUMENT CLASS ──────────────────────────────────────────────────────────
% book class: best for multi-chapter theses (front/main/back matter support)
% 12pt: universal standard (Oxford, Cambridge, APA, Chicago, Cairo Univ)
% a4paper: international standard (Oxford, Cambridge, Egyptian universities)
% oneside: standard for single-sided printing
% openright: chapters start on right-hand (odd) pages — MIT/Oxford standard
\documentclass[12pt, a4paper, oneside, openright]{report}
% For two-sided printing (required by some universities), change to:
% \documentclass[12pt, a4paper, twoside, openright]{report}

% ─── USER CONFIGURATION ─────────────────────────────────────────────────────
\newcommand{\UniversityNameEN}{Faculty of Computers and Information}
\newcommand{\UniversityNameAR}{كلية الحاسبات والمعلومات}
\newcommand{\DepartmentNameEN}{Computer Science Department}
\newcommand{\DepartmentNameAR}{قسم علوم الحاسب}
\newcommand{\UniversityFullEN}{Cairo University}
\newcommand{\ThesisTitle}{Quiz Generation via Knowledge Distillation and Fine-Tuning}
\newcommand{\AuthorName}{Rehab Hamdy}
\newcommand{\SupervisorName}{Prof.\ XXX}
\newcommand{\AcademicYear}{2025--2026}

% ─── PAGE GEOMETRY ───────────────────────────────────────────────────────────
% FCI Cairo: 3cm left (binding), 2.5cm others
% Egyptian Supreme Council: 4cm binding side, 2.5cm others
% MIT/Stanford: 1.5" (38mm) inner for binding
% Compromise: 3.5cm inner (binding), 2.5cm others — satisfies all standards
\usepackage[
  a4paper,
  top=2.5cm,
  bottom=2.5cm,
  inner=3.5cm,                          % binding side (left for English)
  outer=2.5cm,
  headheight=14.5pt,
  headsep=15mm,
  footskip=20mm,
  marginparwidth=0pt,
  marginparsep=0pt
]{geometry}

% ─── LANGUAGE & ENCODING ─────────────────────────────────────────────────────
\usepackage[utf8]{inputenc}              % UTF-8 input
\usepackage[T1]{fontenc}                 % T1 font encoding
\usepackage[english]{babel}              % English hyphenation
% Add arabic option if bilingual content: \usepackage[english,arabic]{babel}

% ─── FONTS ───────────────────────────────────────────────────────────────────
% Times New Roman: universal thesis standard (Oxford, Cambridge, APA, Chicago,
% MIT, Stanford, FCI Cairo, Ain Shams)
\usepackage{fontspec}                    % XeLaTeX / LuaLaTeX font selection
\setmainfont{Times New Roman}            % Body serif font — 12pt standard
\setsansfont{Arial}                      % Sans-serif for captions, headers
\setmonofont{Courier New}[Scale=0.90]    % Monospace for code — 10~11pt equiv.

% Fallback if TNR not available:
% \setmainfont{TeX Gyre Termes}         % Free Times clone
% \setsansfont{TeX Gyre Heros}          % Free Helvetica/Arial clone
% \setmonofont{TeX Gyre Cursor}         % Free Courier clone

% ─── LINE SPACING ────────────────────────────────────────────────────────────
% FCI Cairo: 1.5 line spacing
% Oxford/Cambridge: 1.5 or double
% APA: double
% MIT/Stanford: 1.5 or double
% → 1.5 is the most common for FCI; use \doublespacing if your dept. requires
\usepackage{setspace}
\onehalfspacing

% ─── PARAGRAPH FORMATTING ────────────────────────────────────────────────────
% APA / Chicago / FCI: 0.5" (12.7mm) first-line indent
% APA: NO extra space between paragraphs
\usepackage{indentfirst}                 % indent first paragraph after heading
\setlength{\parindent}{12.7mm}           % 0.5 inch first-line indent
\setlength{\parskip}{0pt}                % no extra vertical space between paras

% ─── SECTIONAL HEADINGS ─────────────────────────────────────────────────────
% Heading hierarchy based on APA 7th + FCI Cairo synthesis:
%   Chapter  → centered, bold, 16pt (FCI standard)
%   Section  → left-aligned, bold, 14pt (FCI standard)
%   Subsection → left-aligned, bold, 12pt (FCI standard)
%   Subsubsection → left-aligned, bold italic, 12pt
%   Paragraph → indented, bold, run-in, period, 12pt
\usepackage{titlesec}
\usepackage{titletoc}

% ---- Chapter (APA L1 equivalent) ----
% FCI: bold, centered, 16pt
% Oxford/Cambridge: centered, bold, 14-16pt
\titleformat{\chapter}[display]
  {\normalfont\centering\bfseries}
  {\MakeUppercase{\chaptertitlename}\ \thechapter}
  {20pt}
  {\fontsize{16pt}{19.2pt}\selectfont\bfseries}
\titlespacing*{\chapter}{0pt}{-10pt}{30pt}

% ---- Section (APA L2 equivalent) ----
% FCI: bold, left-aligned, 14pt
% APA: Flush Left, Bold, Title Case
\titleformat{\section}
  {\normalfont\bfseries\fontsize{14pt}{16.8pt}\selectfont}
  {\thesection}
  {1em}
  {}
\titlespacing*{\section}{0pt}{18pt plus 4pt minus 2pt}{10pt plus 2pt minus 2pt}

% ---- Subsection (APA L3 equivalent) ----
% FCI: bold, left-aligned, 12pt
% APA: Flush Left, Bold Italic, Title Case
\titleformat{\subsection}
  {\normalfont\bfseries\fontsize{12pt}{14.4pt}\selectfont}
  {\thesubsection}
  {1em}
  {}
\titlespacing*{\subsection}{0pt}{14pt plus 4pt minus 2pt}{8pt plus 2pt minus 2pt}

% ---- Subsubsection (APA L4 equivalent) ----
% APA: Indented, Bold, Title Case, Period., Text continues
% Chicago: Flush Left, Italic
\titleformat{\subsubsection}
  {\normalfont\bfseries\itshape}
  {\thesubsubsection}
  {1em}
  {}
\titlespacing*{\subsubsection}{0pt}{12pt plus 4pt minus 2pt}{6pt plus 2pt minus 2pt}

% ---- Paragraph (APA L5 equivalent) ----
% APA: Indented, Bold Italic, Title Case, Period., Text continues
\titleformat{\paragraph}[runin]
  {\normalfont\bfseries\itshape}
  {\theparagraph}
  {1em}
  {}
\titlespacing*{\paragraph}{\parindent}{10pt plus 4pt minus 2pt}{1em plus 2pt minus 2pt}

% ─── TABLE OF CONTENTS FORMATTING ───────────────────────────────────────────
% Oxford-style: clean, professional dot leaders
\titlecontents{chapter}
  [0em]
  {\bfseries\vspace{6pt}}
  {\contentslabel{2em}}
  {\hspace*{2em}}
  {\titlerule*[8pt]{$\cdot$}\contentspage}

\titlecontents{section}
  [2em]
  {}
  {\contentslabel{2.5em}}
  {\hspace*{2.5em}}
  {\titlerule*[8pt]{$\cdot$}\contentspage}

\titlecontents{subsection}
  [4.5em]
  {}
  {\contentslabel{3em}}
  {\hspace*{3em}}
  {\titlerule*[8pt]{$\cdot$}\contentspage}

\titlecontents{subsubsection}
  [7.5em]
  {\small}
  {\contentslabel{3.5em}}
  {\hspace*{3.5em}}
  {\titlerule*[8pt]{$\cdot$}\contentspage}

% ─── NUMBERING DEPTH ────────────────────────────────────────────────────────
\setcounter{secnumdepth}{4}              % number up to subsubsection
\setcounter{tocdepth}{3}                 % show up to subsection in TOC
                                          % (change to 2 for sections only)

% ─── FIGURE & TABLE NUMBERING ───────────────────────────────────────────────
% Per-chapter numbering: Fig 3.1, Table 3.2 (MIT, Oxford, FCI standard)
\usepackage{chngcntr}
\counterwithin{figure}{chapter}
\counterwithin{table}{chapter}
% For continuous numbering (Fig 1, Fig 2, ...), use instead:
% \counterwithout{figure}{chapter}
% \counterwithout{table}{chapter}

% ─── PACKAGES — TABLES ──────────────────────────────────────────────────────
% APA 7th: ONLY 3 horizontal rules — top, header-body separator, bottom
% NO vertical rules. This is the universal academic standard.
\usepackage{booktabs}                    % \toprule, \midrule, \bottomrule
\usepackage{tabularx}                    % flexible-width tables
\usepackage{longtable}                   % tables spanning multiple pages
\usepackage{multirow}                    % merged cells vertically
\usepackage{makecell}                    % line breaks inside cells
\usepackage{array}                       % extended column types
\usepackage{colortbl}                    % colored cells (use sparingly)

% ─── PACKAGES — FIGURES ─────────────────────────────────────────────────────
\usepackage{graphicx}                    % \includegraphics
\usepackage{float}                       % [H] placement specifier
\usepackage{subcaption}                  % sub-figures (a), (b), (c)
\usepackage{wrapfig}                     % text wrapping around figures
\usepackage{tikz}
\usetikzlibrary{
  shapes, arrows, positioning, calc,
  fit, backgrounds, decorations.pathreplacing,
  matrix, shadows, patterns
}

% ─── PACKAGES — CAPTIONS ────────────────────────────────────────────────────
% Universal standard: Table caption ABOVE, Figure caption BELOW
% APA: Table title italic, "Note." below table
\usepackage[
  font={small},
  labelfont={bf},
  labelsep=quad,
  justification=centerlast,
  skip=8pt
]{caption}

% Figure caption: below figure (ALL standards agree)
\captionsetup[figure]{
  font={small},
  labelfont={bf},
  labelsep=quad,
  justification=centerlast,
  skip=6pt,
  position=bottom
}

% Table caption: above table (ALL standards agree)
\captionsetup[table]{
  font={small},
  labelfont={bf},
  labelsep=quad,
  justification=centerlast,
  skip=6pt,
  position=top
}

% Sub-figure caption
\captionsetup[subfigure]{
  font={footnote},
  labelfont={bf},
  labelsep=space,
  justification=centerlast,
  skip=4pt
}

% ─── PACKAGES — HEADERS & FOOTERS ───────────────────────────────────────────
\usepackage{fancyhdr}

% ─── PACKAGES — COLORS ──────────────────────────────────────────────────────
\usepackage{xcolor}
\usepackage[table]{xcolor}

% ─── PACKAGES — MATHEMATICS ─────────────────────────────────────────────────
\usepackage{amsmath}
\usepackage{amssymb}
\usepackage{amsfonts}
\usepackage{mathtools}
\usepackage{bm}

% ─── PACKAGES — ALGORITHMS ──────────────────────────────────────────────────
\usepackage{algorithm}
\usepackage{algpseudocode}
\usepackage{algorithmic}

% ─── PACKAGES — CODE LISTINGS ───────────────────────────────────────────────
\usepackage{listings}

% ─── PACKAGES — HYPERLINKS ──────────────────────────────────────────────────
\usepackage[
  colorlinks=true,
  linkcolor=ThesisBlue,
  citecolor=ThesisGreen,
  urlcolor=ThesisNavy,
  bookmarks=true,
  bookmarksopen=true,
  bookmarksnumbered=true,
  pdfauthor={\AuthorName},
  pdftitle={\ThesisTitle}
]{hyperref}

% ─── PACKAGES — BIBLIOGRAPHY ────────────────────────────────────────────────
% IEEE style: standard for CS departments (FCI, Ain Shams, MIT EECS)
\usepackage[backend=biber, style=ieee]{biblatex}
% \addbibresource{references.bib}        % uncomment and set your .bib file

% ─── PACKAGES — GLOSSARY / ACRONYMS ─────────────────────────────────────────
\usepackage[acronym, toc]{glossaries}
% \makeglossaries

% ─── PACKAGES — MISCELLANEOUS ───────────────────────────────────────────────
\usepackage{enumitem}                    % customized lists
\usepackage{pifont}                      % ding symbols
\usepackage{calc}
\usepackage{ifthen}
\usepackage{afterpage}
\usepackage{emptypage}                   % blank pages have no headers/footers
\usepackage{appendix}
\usepackage{pdflscape}                   % landscape pages
\usepackage{placeins}                    % \FloatBarrier
\usepackage{siunitx}                     % SI units
\usepackage{tcolorbox}                   % colored boxes for callouts
\usepackage{csquotes}                    % proper quoting (required by biblatex)

% ─── COLOR PALETTE ──────────────────────────────────────────────────────────
% Muted academic colors — professional, not distracting
% Inspired by Oxford/Cambridge thesis templates

% ---- Primary ----
\definecolor{ThesisBlue}{RGB}{0, 70, 140}         % headings, links
\definecolor{ThesisNavy}{RGB}{0, 40, 85}           % deep emphasis
\definecolor{ThesisGreen}{RGB}{0, 120, 60}         % citations, positive
\definecolor{ThesisRed}{RGB}{180, 30, 30}           % warnings, negative
\definecolor{ThesisOrange}{RGB}{210, 105, 0}        % highlights

% ---- Table ----
\definecolor{TableHeader}{RGB}{0, 70, 140}          % header row background
\definecolor{TableHeaderText}{RGB}{255, 255, 255}   % header row text
\definecolor{TableRowA}{RGB}{240, 245, 250}         % alternating row A
\definecolor{TableRowB}{RGB}{255, 255, 255}         % alternating row B

% ---- Code ----
\definecolor{CodeBG}{RGB}{248, 248, 248}
\definecolor{CodeFrame}{RGB}{200, 200, 200}
\definecolor{CodeKeyword}{RGB}{0, 0, 180}
\definecolor{CodeString}{RGB}{0, 128, 0}
\definecolor{CodeComment}{RGB}{128, 128, 128}
\definecolor{CodeNumber}{RGB}{160, 100, 0}

% ---- Callout Boxes ----
\definecolor{NoteBoxBG}{RGB}{232, 244, 255}
\definecolor{NoteBoxFrame}{RGB}{0, 70, 140}
\definecolor{WarningBoxBG}{RGB}{255, 243, 224}
\definecolor{WarningBoxFrame}{RGB}{210, 105, 0}
\definecolor{ImportantBoxBG}{RGB}{255, 232, 232}
\definecolor{ImportantBoxFrame}{RGB}{180, 30, 30}

% ─── PAGE STYLE — HEADERS & FOOTERS ────────────────────────────────────────
% Oxford/Cambridge: chapter title left header, page number right
% APA: page number top right (student paper) or running head + page# (pro)
% FCI: page number bottom center is common
% → Using Oxford-style with page number in header

\pagestyle{fancy}
\fancyhf{}

% ---- Header ----
\fancyhead[LE,RO]{\small\thepage}                     % page number
\fancyhead[RE]{\small\itshape\nouppercase{\leftmark}}  % chapter title (even)
\fancyhead[LO]{}                                       % leave blank (odd)
\renewcommand{\headrulewidth}{0.4pt}

% ---- Footer ----
\fancyfoot[C]{}
\fancyfoot[L]{\tiny\textcolor{gray}{\UniversityNameEN}}

% ---- Chapter-opening pages: no header, page number bottom center ----
\fancypagestyle{plain}{%
  \fancyhf{}%
  \fancyfoot[C]{\small\thepage}%
  \renewcommand{\headrulewidth}{0pt}%
}

% ─── TABLE FORMATTING ───────────────────────────────────────────────────────
% APA 7th (universally accepted):
%   - ONLY 3 horizontal rules: \toprule, \midrule, \bottomrule
%   - NO vertical rules
%   - Caption ABOVE table
%   - "Note." line below table (optional)
% Table text may be smaller (10pt) than body (12pt) — APA, Chicago

\renewcommand{\arraystretch}{1.30}       % row height (APA recommends 1.25-1.35)
\tabcolsep=8pt                           % column separation

% ---- Custom column types ----
\newcolumntype{L}[1]{>{\raggedright\arraybackslash}p{#1}}
\newcolumntype{C}[1]{>{\centering\arraybackslash}p{#1}}
\newcolumntype{R}[1]{>{\raggedleft\arraybackslash}p{#1}}
\newcolumntype{M}[1]{>{\centering\arraybackslash}m{#1}}

% ---- Styled table header command ----
% Use within booktabs tables for professional APA-style headers
\newcommand{\thead}[1]{\textbf{#1}}

% ---- Table Note command (APA style) ----
% Usage: \tablenote{Source: Author's calculations. n = 200.}
\newcommand{\tablenote}[1]{%
  \vspace{2pt}%
  {\small\textit{Note.}\quad #1}%
}

% ---- Professional table environment ----
% Usage:
%   \begin{proftable}{caption}{label}{col spec}
%     \toprule
%     \thead{Col 1} & \thead{Col 2} \\
%     \midrule
%     data & data \\
%     \bottomrule
%   \end{proftable}
\newenvironment{proftable}[3]{%
  \begin{table}[H]%
  \centering%
  \caption{#1}%
  \label{#2}%
  \small%
  \begin{tabular}{#3}%
}{%
  \end{tabular}%
  \end{table}%
}

% ---- Long table environment (multi-page) ----
% Same rules apply: booktabs only, no vertical rules
\newenvironment{proflongtable}[3]{%
  \begin{longtable}{#3}%
  \caption{#1}\label{#2}\\%
}{%
  \end{longtable}%
}

% ─── FIGURE PLACEMENT ───────────────────────────────────────────────────────
% APA: figures embedded near first mention, not floating far away
\renewcommand{\topfraction}{0.92}
\renewcommand{\bottomfraction}{0.85}
\renewcommand{\textfraction}{0.08}
\renewcommand{\floatpagefraction}{0.90}

% ─── LIST FORMATTING ────────────────────────────────────────────────────────
% APA: bullet lists allowed, no extra spacing between items
% Chicago: consistent indentation

\setitemize{
  leftmargin=2em,
  itemsep=4pt,
  parsep=2pt,
  topsep=6pt
}

\setenumerate{
  leftmargin=2em,
  itemsep=4pt,
  parsep=2pt,
  topsep=6pt,
  label=\arabic*.
}

\setdescription{
  leftmargin=2em,
  itemsep=4pt,
  parsep=2pt,
  topsep=6pt,
  font=\normalfont\bfseries
}

% ─── CODE LISTING FORMATTING ────────────────────────────────────────────────
% FCI Cairo: Courier New 10pt for code
\lstdefinestyle{ThesisCode}{
  backgroundcolor=\color{CodeBG},
  basicstyle=\ttfamily\footnotesize,
  keywordstyle=\color{CodeKeyword}\bfseries,
  stringstyle=\color{CodeString},
  commentstyle=\color{CodeComment}\itshape,
  numberstyle=\tiny\color{CodeNumber},
  numbers=left,
  numbersep=8pt,
  frame=single,
  rulecolor=\color{CodeFrame},
  framesep=5pt,
  xleftmargin=15pt,
  xrightmargin=5pt,
  breaklines=true,
  showstringspaces=false,
  tabsize=4,
  captionpos=b,
  columns=flexible,
  keepspaces=true,
  escapeinside={(*@}{@*)}
}
\lstset{style=ThesisCode}

\lstdefinestyle{PythonCode}{style=ThesisCode, language=Python}
\lstdefinestyle{BashCode}{style=ThesisCode, language=bash}

\newcommand{\code}[1]{\texttt{\small #1}}

% ─── ALGORITHM / PSEUDOCODE FORMATTING ──────────────────────────────────────
\algrenewcommand\algorithmicrequire{\textbf{Input:}}
\algrenewcommand\algorithmicensure{\textbf{Output:}}
\floatname{algorithm}{Algorithm}
\renewcommand{\algorithmiccomment}[1]{\hfill$\triangleright$ \textit{#1}}

% ─── MATHEMATICS FORMATTING ─────────────────────────────────────────────────
% Per-chapter equation numbering: (3.1), (3.2) — MIT, Oxford, Springer standard
\DeclareMathOperator*{\argmin}{arg\,min}
\DeclareMathOperator*{\argmax}{arg\,max}
\DeclareMathOperator{\KL}{KL}
\DeclareMathOperator{\softmax}{softmax}
\DeclareMathOperator{\sigmoid}{sigmoid}
\DeclareMathOperator{\ReLU}{ReLU}

% Math shortcuts
\newcommand{\R}{\mathbb{R}}
\newcommand{\E}{\mathbb{E}}
\newcommand{\norm}[1]{\left\lVert #1 \right\rVert}
\newcommand{\abs}[1]{\left\lvert #1 \right\rvert}
\newcommand{\inner}[2]{\left\langle #1,\, #2 \right\rangle}
\newcommand{\condbar}{\,\vert\,}
\newcommand{\bigO}{\mathcal{O}}

\numberwithin{equation}{chapter}

% ─── CALLOUT BOXES ──────────────────────────────────────────────────────────
% Use sparingly — academic theses prefer inline text over boxes
% But useful for: definitions, key findings, warnings

\newtcolorbox{notebox}[1][Note]{%
  colback=NoteBoxBG,
  colframe=NoteBoxFrame,
  fonttitle=\bfseries\sffamily,
  title=#1,
  boxrule=0.8pt,
  arc=2pt,
  left=8pt, right=8pt, top=4pt, bottom=4pt,
  before skip=10pt, after skip=10pt
}

\newtcolorbox{warningbox}[1][Warning]{%
  colback=WarningBoxBG,
  colframe=WarningBoxFrame,
  fonttitle=\bfseries\sffamily,
  title=#1,
  boxrule=0.8pt,
  arc=2pt,
  left=8pt, right=8pt, top=4pt, bottom=4pt,
  before skip=10pt, after skip=10pt
}

\newtcolorbox{importantbox}[1][Important]{%
  colback=ImportantBoxBG,
  colframe=ImportantBoxFrame,
  fonttitle=\bfseries\sffamily,
  title=#1,
  boxrule=0.8pt,
  arc=2pt,
  left=8pt, right=8pt, top=4pt, bottom=4pt,
  before skip=10pt, after skip=10pt
}

\newtcolorbox{definitionbox}[1][Definition]{%
  colback=NoteBoxBG,
  colframe=ThesisBlue,
  fonttitle=\bfseries\sffamily\color{white},
  coltitle=white,
  colbacktitle=ThesisBlue,
  title=#1,
  boxrule=0.8pt,
  arc=2pt,
  left=8pt, right=8pt, top=4pt, bottom=4pt,
  before skip=10pt, after skip=10pt
}

% ─── CROSS-REFERENCE COMMANDS ───────────────────────────────────────────────
% Friendly references — universal academic style
\newcommand{\reffig}[1]{Figure~\ref{#1}}
\newcommand{\reftab}[1]{Table~\ref{#1}}
\newcommand{\refeq}[1]{Equation~\ref{#1}}
\newcommand{\refch}[1]{Chapter~\ref{#1}}
\newcommand{\refsec}[1]{Section~\ref{#1}}
\newcommand{\refalg}[1]{Algorithm~\ref{#1}}
\newcommand{\refapp}[1]{Appendix~\ref{#1}}
\newcommand{\reflst}[1]{Listing~\ref{#1}}
\newcommand{\refpage}[1]{page~\pageref{#1}}

% ─── UTILITY COMMANDS ───────────────────────────────────────────────────────
\newcommand{\pct}[1]{#1\%}
\newcommand{\metric}[2]{\textbf{#1}\,=\,#2}

% ─── COVER PAGE ─────────────────────────────────────────────────────────────
% FCI Cairo style: University name, Faculty, Department, Title, Author, Supervisor
\newcommand{\MakeCoverPage}{%
  \begin{titlepage}
    \centering
    \vspace*{0.5cm}

    {\scshape\LARGE \UniversityFullEN \par}
    \vspace{0.3cm}
    {\scshape\Large \UniversityNameEN \par}
    \vspace{0.2cm}
    {\scshape\large \DepartmentNameEN \par}

    \vspace{2cm}

    {\fontsize{18pt}{21.6pt}\selectfont\bfseries \ThesisTitle \par}

    \vspace{1.5cm}

    {\Large Graduation Project Thesis\par}

    \vspace{2cm}

    \begin{minipage}{0.45\textwidth}
      \begin{flushleft}\large
        \textbf{Submitted by:}\\
        \AuthorName
      \end{flushleft}
    \end{minipage}
    \hfill
    \begin{minipage}{0.45\textwidth}
      \begin{flushright}\large
        \textbf{Supervised by:}\\
        \SupervisorName
      \end{flushright}
    \end{minipage}

    \vfill

    {\large \AcademicYear\par}

  \end{titlepage}
}

% ─── ABSTRACT PAGE ──────────────────────────────────────────────────────────
% FCI: Arabic abstract first, then English (or vice versa per department)
\newcommand{\MakeAbstractPage}[2]{%
  \chapter*{Abstract}
  \addcontentsline{toc}{chapter}{Abstract}
  #1
  \vfill
  \ifthenelse{\equal{#2}{}}{}{%
    \newpage
    \chapter*{الملخص}
    \addcontentsline{toc}{chapter}{الملخص}
    #2
  }
}

% ─── ACKNOWLEDGMENTS ────────────────────────────────────────────────────────
\newcommand{\MakeAcknowledgmentsPage}[1]{%
  \chapter*{Acknowledgments}
  \addcontentsline{toc}{chapter}{Acknowledgments}
  #1
}

% ─── DEDICATION ─────────────────────────────────────────────────────────────
\newcommand{\MakeDedicationPage}[1]{%
  \vspace*{\fill}
  \begin{center}
    \textit{\Large #1}
  \end{center}
  \vspace*{\fill}
}

% ─── APPENDIX FORMATTING ────────────────────────────────────────────────────
% APA: labeled Appendix A, Appendix B; each on new page
% FCI: similar structure
\renewcommand{\appendixname}{Appendix}

% ─── BIBLIOGRAPHY FORMATTING ────────────────────────────────────────────────
% IEEE style for CS departments (FCI, Ain Shams, MIT EECS)
\defbibheading{bibliography}{%
  \chapter*{Bibliography}%
  \addcontentsline{toc}{chapter}{Bibliography}%
}

% ─── FRONT MATTER / MAIN MATTER ─────────────────────────────────────────────
% MIT/Oxford/Cambridge standard:
%   Front matter: Roman numerals (i, ii, iii, ...)
%   Main matter:  Arabic numerals (1, 2, 3, ...)
%   Chapter 1 starts at page 1
\makeatletter
\renewcommand\frontmatter{%
  \cleardoublepage
  \pagenumbering{roman}
}

\renewcommand\mainmatter{%
  \cleardoublepage
  \pagenumbering{arabic}
}

\renewcommand\backmatter{%
  \cleardoublepage
  \setcounter{chapter}{0}%
}
\makeatother

% ─── DRAFT MARKERS ──────────────────────────────────────────────────────────
\newif\ifdraft
\drafttrue                              % set to \draftfalse for final version

\ifdraft
  \newcommand{\todo}[1]{\textcolor{ThesisRed}{\bfseries[TODO: #1]}}
  \newcommand{\FIXME}[1]{\textcolor{ThesisRed}{\bfseries[FIXME: #1]}}
  \newcommand{\CITE}[1]{\textcolor{ThesisOrange}{\bfseries[CITE: #1]}}
\else
  \newcommand{\todo}[1]{}
  \newcommand{\FIXME}[1]{}
  \newcommand{\CITE}[1]{\cite{#1}}
\fi

% ─── ORPHAN / WIDOW PREVENTION ──────────────────────────────────────────────
% Chicago/Oxford: avoid orphan/widow lines
\clubpenalty=10000
\widowpenalty=10000
\emergencystretch=1em
\tolerance=200
\setlength{\emergencystretch}{3em}

% ─── HYPHENATION EXCEPTIONS ─────────────────────────────────────────────────
\hyphenation{
  fine-tun-ing
  know-ledge
  distil-la-tion
  eval-ua-tion
  pre-train-ed
  multi-ple
  parameter
  hyper-parameter
  archi-tec-ture
}

% ─── SI UNIT FORMATTING ─────────────────────────────────────────────────────
\sisetup{
  detect-all,
  group-separator={,},
  group-minimum-digits=4,
  per-mode=symbol
}

% =============================================================================
% FORMATTING DECISIONS LOG — Based on Cross-Source Research
% =============================================================================
%
% | Decision              | Chosen Value       | Sources Supporting It          |
% |-----------------------|--------------------|--------------------------------|
% | Page size             | A4                 | Oxford, Cambridge, Cairo, FCI  |
% | Inner margin          | 3.5cm              | FCI (3cm), MIT (1.5"), avg     |
% | Outer margin          | 2.5cm              | FCI, Egyptian Council          |
% | Body font             | TNR 12pt           | ALL sources agree              |
% | Line spacing          | 1.5                | FCI, Cambridge, Oxford         |
% | Paragraph indent      | 12.7mm (0.5")      | APA, Chicago, FCI              |
% | Chapter heading       | 16pt bold centered | FCI, Oxford, Cambridge         |
% | Section heading       | 14pt bold left     | FCI standard                   |
% | Subsection heading    | 12pt bold left     | FCI, APA L3                    |
% | Table rules           | 3 only (booktabs)  | APA 7th, ALL academic stds     |
% | Table caption         | Above table        | APA, IEEE, ACM, ALL sources    |
% | Figure caption        | Below figure       | APA, IEEE, ACM, ALL sources    |
% | Equation numbering    | Per chapter        | MIT, Oxford, Springer          |
% | Figure numbering      | Per chapter        | MIT, Oxford, FCI               |
% | Bibliography style    | IEEE               | FCI CS dept, ACM, MIT EECS     |
% | Front matter pages    | Roman numerals     | MIT, Oxford, Cambridge, FCI    |
% | Main matter pages     | Arabic numerals    | ALL sources agree              |
% | TOC depth             | 3 levels           | Oxford, Cambridge              |
% | Secnum depth          | 4 levels           | Standard for theses            |
%
% =============================================================================

% =============================================================================
% END OF FORMATTING FILE v2
% =============================================================================
